\textbf{(D) 定理 2 和 3 的证明。}

李代数同态 \(\beta \circ \alpha: \mathrm{L}_{\mathbf{Z}}(\mathrm{X}) \to \mathrm{A}(\mathrm{X})\) 限制于 \(\mathbf{X}\) 后，由 (3) 可知是恒等映射，因而是典范嵌入（§ 3，第 1 号）。所以 \(\alpha\) 是单射，并且由 (B) 是双射；这证明了定理 3。由于 \(\beta \circ \alpha = f \circ \varepsilon \circ \alpha\) 是单射且 \(\alpha\) 是双射，\(\varepsilon\) 是单射。对所有 \(n \geqslant 1\)，

\[
\varepsilon_ {n} \colon \mathrm{C} ^ {n} / \mathrm{C} ^ {n + 1} \rightarrow \mathrm{F} ^ {n} / \mathrm{F} ^ {n + 1}
\]

是单射，因而

\[
\mathrm{C} ^ {n} \cap \mathrm{F} ^ {n + 1} = \mathrm{C} ^ {n + 1}.
\]

\(C^{1} = F = F^{1}\)；若 \(C^{n} = F^{n}\)，则 \(C^{n} \cap F^{n+1} = F^{n+1}\)，从而 \(C^{n+1} = F^{n+1}\)，这通过对 \(n \geqslant 1\) 作归纳证明了定理 2。

推论。

\[
\bigcap_ {n \geqslant 1} \mathrm{C} ^ {n} \mathrm{F} (\mathrm{X}) = \{e \}.
\]

取 K = Z 且 r = 0，应用定理 2，

\[
\bigcap_ {n \geqslant 1} \mathrm{C} ^ {n} \mathrm{F} (\mathrm{X}) = \bigcap_ {n \geqslant 1} \bar {g} (1 + \hat {\mathrm{A}} _ {n} (\mathrm{X})) = \bar {g} \left(\bigcap_ {n \geqslant 1} (1 + \hat {\mathrm{A}} _ {n} (\mathrm{X}))\right) = \bar {g} (1) = \{e \}.
\]

注。设 H 是关于 X 的 Hall 集（§2，第 10 号）。设 M 是在 F(X) 上由复合律 \((x,y)\mapsto (x,y) = x^{-1}y^{-1}xy\) 定义的二元代数系统，并设 \(\phi\) 是从 M(X) 到 M 的同态，其在 X 上的限制为恒等映射。称 \(\phi (\mathrm{H})\) 的元素为与 Hall 集 H 相关的 F(X) 的基本换位子。对每个整数 \(n\geqslant 1\)，令 \(\mathrm{H}_{\mathfrak{n}}\) 为 H 的子集，它由长度为 \(n\) 的元素组成；我们知道（§2，第 11 号，定理 1），\(\mathrm{H}_{\mathfrak{n}}\) 到 \(\mathbf{L}_{\mathbf{Z}}(\mathbf{X})\) 的典范映射是阿贝尔群 \(\mathbf{L}_{\mathbf{Z}}^{n} (\mathbf{X})\) 的一组基。此外，\(\phi (\mathrm{H_n})\subset \mathrm{C^n}\)；对于所有 \(m\in \mathrm{H}_{\mathfrak{n}}\)，令 \(\phi_{\mathfrak{n}}(m)\) 表示 \(\mathrm{C}^{n + 1}\) 模 \(\phi (m)\in \mathrm{C}^n\) 的陪集。于是定理 3 表明，\(\phi_{\mathfrak{n}}\) 是从 \(\mathrm{H}_{\mathfrak{n}}\) 到阿贝尔群 \(\mathrm{C}^n /\mathrm{C}^{n + 1}\) 的一组基的双射。立即得到，对于所有 \(w\in F(\mathbf{X})\) 和所有 \(i\geqslant 1\)，存在 \(\alpha_{i}\) 中唯一的元素 \(\mathbf{Z}^{(\mathrm{H}_{i})}\)，使得对于 \(n\geqslant 1\)，

\[
w = \prod_ {i = 1} ^ {n} \prod_ {m \in H _ {i}} \phi (m) ^ {\alpha_ {i} (m)} \mod C ^ {n + 1},\tag{4}
\]

其中，乘积按照 H 上给出的全序计算。

例。设 X 是含有两个元素 x、y 的集合，并令 \(H_{1} = \{x, y\}\)、\(H_{2} = \{xy\}\)。因而 F(X) 的每个元素 w 都可以写成

\[
w \equiv x ^ {a} y ^ {b} (x, y) ^ {c} \bmod . \mathrm{C} ^ {3} \quad \text { with } a, b, c \text { in } \mathbf {Z}.
\]

对于 \(w = (xy)^n\)，有 \(a = b = n\) 且 \(c = n(1 - n)/2\)（参见习题 9），从而

\[
(x y) ^ {n} \equiv x ^ {n} y ^ {n} (x, y) ^ {n (1 - n) / 2} \mod C ^ {3}.
\]

\subsection{自由群的 p-滤子}

在本号中，\(p\) 表示一个素数，并假设 \(\mathbf{K} = \mathbf{F}_p\)。令 \(\gamma\) 为从 \(\mathrm{F}(\mathbf{X})\) 到 \(\Gamma(\mathbf{X})\) 的同态，它由对 \(\gamma(x) = 1 + x\) 中的 \(x\) 规定 \(\mathbf{X}\)；记 \(\mathrm{F}_n^{(p)}(\mathbf{X}) = \gamma(1 + \hat{\Lambda}_n(\mathbf{X}))\)。序列 \((\mathrm{F}_n^{(p)}(\mathbf{X}))_{n\geq 1}\) 是 \(\mathrm{F}(\mathbf{X})\) 上的一个整中心滤子，由于 \(\gamma\) 是单射（第 3 号，定理 1），它是分离的。它称为 \(p\) 上的 \(\mathrm{F}(\mathbf{X})\)-滤子。

命题 2。假设 X 是有限集。对每个整数 \(n \geqslant 1\)，群 \(F(X) / F_n^{(p)}(X)\) 是幂零类 \(p\) 的有限 \(\leqslant n\)-群。

按 \(n\) 归纳论证，只需证明对所有 \(\mathrm{F}_{n}^{(p)}(\mathbf{X}) / \mathrm{F}_{n + 1}^{(p)}(\mathbf{X})\)，\(p\) 是有限交换 \(n \geqslant 1\)-群。对所有 \(w \in \mathrm{F}_{n}^{(p)}(\mathbf{X})\)，元素 \(\gamma(w) - 1\) 属于 \(\hat{\mathbf{A}}(\mathbf{X})\)，且其阶 \(\geqslant n\)；记 \(\delta_n(w)\) 为 \(\gamma(w) - 1\) 的 n 次齐次分量。映射 \(n\) 是一个同态，其核为 \(\delta_n: \mathrm{F}_{n}^{(p)}(\mathbf{X}) \to \mathrm{A}^n (\mathbf{X})\)（§ 4，第 5 号，命题 3），因而 \(\mathrm{F}_{n + 1}^{(p)}(\mathbf{X})\) 同构于 \(\mathrm{F}_{n}^{(p)}(\mathbf{X}) / \mathrm{F}_{n + 1}^{(p)}(\mathbf{X})\) 的一个子群。由于 \(\mathrm{A}^n (\mathbf{X})\) 是有限集，\(\mathbf{X}\) 是 \(\mathrm{A}^n (\mathbf{X})\) 上的有限维向量空间，因而是有限交换 \(\mathbf{F}_p\)-群，所以 \(p\) 也是有限交换 \(\mathrm{F}_{n}^{(p)}(\mathbf{X}) / \mathrm{F}_{n + 1}^{(p)}(\mathbf{X})\)-群。

命题 3。对于 \(w \neq 1\) 中所有满足 \(\mathrm{F}(\mathbf{X})\) 的 w，存在一个有限 \(p\)-群 \(\mathbf{G}\) 以及一个从 \(f\) 到 \(\mathrm{F}(\mathbf{X})\) 的同态 \(\mathbf{G}\)，使得 \(f(w) \neq 1\)。

存在 \(x_{1},\ldots ,x_{r}\) 中的元素 \(\mathbf{X}\) 和整数 \(n_1,\dots ,n_r\)，使得 \(w = x_1^{n_1}\dots x_r^{n_r}\)。令 \(\mathrm{Y} = \{x_1,\dots,x_r\}\)。\(\mathrm{Y}\) 到 \(\mathbf{X}\) 的典范嵌入延拓为同态 \(\alpha : \mathrm{F}(\mathrm{Y})\to \mathrm{F}(\mathrm{X})\)；另一方面，令 \(\beta\) 为从 \(\mathrm{F}(\mathrm{X})\) 到 \(\mathrm{F}(\mathrm{Y})\) 的同态，其在 \(\mathbf{Y}\) 上的限制为恒等映射，并将 \(\mathbf{X} - \mathbf{Y}\) 映到 \(\{1\}\)。于是对 \(\beta (\alpha (y)) = y\)，有 \(y\in \mathbf{Y}\)，从而 \(\beta \circ \alpha\) 是 \(\mathrm{F}(\mathrm{Y})\) 的恒等自同构。显然存在 \(w^{\prime}\) 中的 \(\mathrm{F}(\mathrm{Y})\)，使得 \(w = \alpha (w^{\prime})\)；于是 \(\beta (w) = w^{\prime}\neq 1\)；现在 \(\bigcap_{n\geqslant 1}\mathbf{F}_n^{(p)}(\mathbf{Y}) = \{1\}\)，因此存在整数 \(n\geqslant 1\)，使得 \(\beta (w)\notin \mathbf{F}_{\mathrm{n}}^{(p)}(\mathbf{Y})\)。由命题 2，群 \(\mathbf{G} = \mathbf{F}(\mathbf{Y}) / \mathbf{F}_{\mathrm{n}}^{(p)}(\mathbf{Y})\) 是有限 \(p\)-群。若 \(f\) 是 \(\beta\) 与从 \(\mathrm{F}(\mathrm{Y})\) 到 \(\mathbf{G}\) 的典范同态的复合，则 \(f(w)\neq 1\)。

推论。F(X) 中所有有限指标正规子群的交为 \{1\}。

\section{豪斯多夫级数}

本段中假设 K 是特征为 0 的域。

\subsection{滤代数中的指数与对数}

设 \(\mathbf{A}\) 是一个含幺结合代数，在实滤子 \((\mathbf{A}_{\alpha})\) 下是豪斯多夫且完备的。记 \(m = A_0^+ = \bigcup_{\alpha > 0} A_\alpha\)。

对于 \(x \in \mathfrak{m}\)，族 \((x^n / n!)_{n \in \mathbf{N}}\) 是可求和的。记

\[
e ^ {x} = \exp x = \sum_ {n \geqslant 0} x ^ {n} / n!.\tag{1}
\]

于是 \(\exp (x)\in 1 + \mathfrak{m}\)，映射 \(\exp \colon \mathfrak{m}\to 1 + \mathfrak{m}\) 称为 A 的指数映射。

对于所有 \(y \in 1 + \mathfrak{m}\)，族 \(((-1)^{n-1}(y - 1)^n / n)_{n \geqslant 1}\) 是可求和的。记

\[
\log y = \sum_ {n \geqslant 1} (- 1) ^ {n - 1} (y - 1) ^ {n} / n.\tag{2}
\]

于是 \(\log y \in m\)，映射 \(\log: 1 + m \to m\) 称为 A 的对数映射。

命题 1。指数映射是从 \(\mathfrak{m}\) 到 \(1 + \mathfrak{m}\) 的同胚，而对数映射是其逆同胚。

对于 \(x \in \mathrm{A}_{\alpha}, \frac{x^n}{n!} \in \mathrm{A}_{n\alpha}\)，有 \(\mathrm{A}_{\alpha}\)。因此，定义指数的级数在每个 \(\alpha > 0\) 的集合 \(\mathrm{A}_{\alpha}\) 上一致收敛；由于 \(m\) 在 m 中是开集且 \(m = \bigcup_{\alpha > 0} \mathrm{A}_{\alpha}\)，指数映射连续。类似地可证明对数映射连续。

设 \(e\) 和 \(l\) 是没有常数项的形式幂级数

\[
e (\mathrm{X}) = \sum_ {n \geqslant 1} \frac {\mathrm{X} ^ {n}}{n !}, \quad l (\mathrm{X}) = \sum_ {n \geqslant 1} (- 1) ^ {n - 1} \mathrm{X} ^ {n} / n.
\]

我们知道（《代数》，第 IV 章，§ 6，第 9 号），在 \(e(l(\mathbf{X})) = l(e(\mathbf{X})) = \mathbf{X}\) 上有 \(\hat{\mathrm{A}}(\{\mathbf{X}\}) = \mathrm{K}[[\mathbf{X}]]\)。由代入（§ 5，第 1 号），得到

\[
e (l (x)) = l (e (x)) = x
\]

对于 \(x\in \mathfrak{m}\)；由于

\[
\exp x = e (x) + 1, \quad \log (1 + x) = l (x)
\]

于是立即得到

\[
\log \exp x = x, \quad \exp \log (1 + x) = 1 + x
\]

对于 \(x\) 中的 \(\mathfrak{m}\)，从而命题得证。

注。（1）若 \(x \in \mathfrak{m}\)、\(y \in \mathfrak{m}\) 且 \(x\) 与 \(y\) 可交换，则

\[
\exp (x + y) = \exp (x) \exp (y),
\]

因为族 \(\left(\frac{x^i}{i!}\cdot \frac{y^j}{j!}\right)_{i,j\in \mathbf{N}}\) 是可求和的（参见《代数》，第 IV 章，§ 6，第 9 号，命题 11）。

\begin{enumerate}
\def\labelenumi{(\arabic{enumi})}
\setcounter{enumi}{1}
\item
  由于级数 \(e\) 和 \(l\) 不含常数项且 \(A_{\alpha}\) 是 \(A\) 的闭理想，有 \(\exp A_{\alpha} \subset 1 + A_{\alpha}\) 和 \(\log (1 + A_{\alpha}) \subset A_{\alpha}\)，从而当 \(\exp A_{\alpha} = 1 + A_{\alpha}\) 时，\(\log (1 + A_{\alpha}) = A_{\alpha}\) 且 \(\alpha > 0\)。
\item
  设 B 是一个完备豪斯多夫滤含幺结合代数，且 \(n = \bigcup_{\alpha > 0} B_{\alpha}\)。设 f 是从 A 到 B 的连续含幺同态，满足 \(f(\mathfrak{m}) \subset \mathfrak{n}\)。则对于 \(f(\exp x) = \exp f(x)\)，有 \(x \in \mathfrak{m}\)；对于 \(f(\log y) = \log f(y)\)，有 \(y \in 1 + \mathfrak{m}\)；例如我们证明第一个公式：
\end{enumerate}

\[
f (\exp x) = \sum_ {n \geqslant 0} f (x ^ {n}) / n! = \sum_ {n \geqslant 0} f (x) ^ {n} / n! = \exp f (x).
\]

\begin{enumerate}
\def\labelenumi{(\arabic{enumi})}
\setcounter{enumi}{3}
\tightlist
\item
  设 E 是一个含幺结合代数。若 \(a\) 是 E 的幂零元素，则族 \(\left(\frac{a^n}{n!}\right)_{n\in \mathbf{N}}\) 具有有限支集，记 \(\exp a = \sum_{n\geqslant 0}a^n /n!\)。若 \(b\) 是幂零的，则元素 \(b - 1\) 称为幂单元；于是记
\end{enumerate}

\[
\log b = \sum_ {n \geqslant 1} (- 1) ^ {n - 1} (b - 1) ^ {n} / n.
\]

由关系 \(e(l(\mathbf{X})) = l(e(\mathbf{X})) = \mathbf{X}\) 可推出，映射 \(a \mapsto \exp a\) 是从 E 的幂零元素集到 E 的幂单元元素集的双射，而 \(b \mapsto \log b\) 是其逆映射。

\subsection{豪斯多夫群}

设 \(\mathbf{X}\) 是一个集合。使用 § 5 第 1、2 号的记号。自由李代数 \(\mathrm{L}(\mathbf{X})\) 被认作其在 \(\mathrm{A}(\mathbf{X})\) 中的典范像（§ 3，第 1 号，定理 1）。记 \(\hat{\mathrm{L}}(\mathbf{X})\) 为 \(\mathrm{L}(\mathbf{X})\) 在 \(\hat{\mathrm{A}}(\mathbf{X})\) 中的闭包，即 \(\hat{\mathrm{A}}(\mathbf{X})\) 中形如 \(a = \sum_{n\geqslant 1}a_n\) 且满足对所有 \(a_{n}\in \mathrm{L}^{n}(\mathrm{X})\) 都有 \(n\geqslant 0\) 的元素所成的集合；这是 \(\hat{\mathrm{A}}(\mathbf{X})\) 的滤李子代数。

定理 1。将 \(\hat{\mathbf{A}}(\mathbf{X})\) 的指数映射限制在 \(\hat{\mathbf{L}}(\mathbf{X})\) 上，得到从 \(\hat{\mathbf{L}}(\mathbf{X})\) 到 Magnus 群 \(\Gamma(\mathbf{X})\) 的一个闭子群的双射。

记 \(\mathrm{A}(\mathrm{X}) = \mathrm{A}\)、\(\mathrm{A}^{\mathrm{n}}(\mathrm{X}) = \mathrm{A}^{\mathrm{n}}\)、\(\hat{\mathrm{A}}(\mathrm{X}) = \hat{\mathrm{A}}\)、\(\mathrm{L}^{\mathrm{n}}(\mathrm{X}) = \mathrm{L}^{\mathrm{n}}\)、\(\hat{\mathrm{L}}(\mathrm{X}) = \hat{\mathrm{L}}\)、\(\Gamma(\mathrm{X}) = \Gamma\)。令 \(\mathrm{B}\) 为代数 \(\mathrm{A} \otimes \mathrm{A}\)，其带有由 \(\mathbf{N}\) 定义的 \(\mathrm{B}^{\mathrm{n}} = \sum_{i+j=n} \mathrm{A}^{i} \otimes \mathrm{A}^{j}\) 型分次。令 \(\hat{\mathrm{B}} = \prod_{n > 0} \mathrm{B}^{n}\) 为相应的完备滤代数（《交换代数》，第 III 章，§2，第 12 号，例 1）。§3 第 1 号定理 1 的推论 1 中定义的余乘法 \(c: \mathrm{A} \to \mathrm{A} \otimes \mathrm{A}\) 是 0 次齐次的，因而连续延拓为同态 \(\hat c: \hat{\mathrm{A}} \to \hat{\mathrm{B}}\)，其给出为

\[
\hat {c} \left(\sum_ {n \geqslant 0} a _ {n}\right) = \sum_ {n \geqslant 0} c (a _ {n}) \quad \text { for } a _ {n} \in \mathrm{A} ^ {n}.
\]

我们还定义从 \(\delta'\) 到 \(\delta''\) 的连续同态 \(\hat{\mathbf{A}}\) 和 \(\hat{\mathbf{B}}\)，规定为

\[
\delta^ {\prime} \left(\sum_ {n \geqslant 0} a _ {n}\right) = \sum_ {n \geqslant 0} (a _ {n} \otimes 1), \quad \delta^ {\prime \prime} \left(\sum_ {n \geqslant 0} a _ {n}\right) = \sum_ {n \geqslant 0} (1 \otimes a _ {n}) \quad \text { for } a _ {n} \in \mathrm{A} ^ {n}.
\]

根据 §3 第 1 号定理 1 的推论 2，\(\mathbf{L}^n\) 是满足 \(a_{n}\in \mathbf{A}^{n}\) 的 \(c(a_{n}) = a_{n}\otimes 1 + 1\otimes a_{n}\) 所成的集合。由此，\(\hat{\mathbf{L}}\) 是满足下式的 \(a\in \hat{\mathbf{A}}\) 所成的集合：

\[
\hat {c} (a) = \delta^ {\prime} (a) + \delta^ {\prime \prime} (a).\tag{3}
\]

令 \(\Delta\) 为 \(b\in \hat{\mathbf{A}}\) 中常数项等于 1 且满足关系

\[
\hat {c} (b) = \delta^ {\prime} (b). \delta^ {\prime \prime} (b),\tag{4}
\]

换言之，是满足 \(b = \sum_{n\geqslant 0}b_n\)、\(b_{n}\in \mathbf{A}^{n}\)（对所有 \(n\geqslant 0\)）、\(b_0 = 1\) 且满足 \(c(b_{n}) = \sum_{t + j = n}b_{t}\otimes b_{j}\)（对 \(n\geqslant 0\)）的 b 所成的集合。后一个刻画表明 \(\Delta\) 是 \(\Gamma\) 的闭子集；由于 \(\hat c,\delta^{\prime}\) 和 \(\delta ''\) 是环同态，且 \(\delta'(\hat{\mathbf A})\) 的每个元素都与 \(\delta''(\hat{\mathbf A})\) 的每个元素可交换，在 \(\Gamma\) 上，映射 \(c\) 和 \(\delta ' \delta''\) 的限制是群同态，而 \(\Delta\) 是 \(\Gamma\) 的子群。

由第 1 号命题 1，\(\hat{\mathbf{A}}\) 的指数映射是从无常数项的 \(\hat{\mathbf{A}}^{+}\) 元素集 \(\hat{\mathbf{A}}\) 到 \(\Gamma\) 的双射。令 \(a \in \hat{\mathbf{A}}^{+}\) 且 \(b = \exp a\)。由于 \(\hat c\) 是连续环同态，

\[
\hat {c} (b) = \hat {c} \left(\sum_ {n \geqslant 0} a ^ {n} / n!\right) = \sum_ {n \geqslant 0} \hat {c} (a) ^ {n} / n! = \exp \hat {c} (a).
\]

关系

\[
\delta^ {\prime} (b) = \exp \delta^ {\prime} (a), \quad \delta^ {\prime \prime} (b) = \exp \delta^ {\prime \prime} (a)
\]

类似可证；又由于 \(\delta'(a)\) 与 \(\delta''(a)\) 可交换（第 1 号，注 1），

\[
\delta^ {\prime} (b) \delta^ {\prime \prime} (b) = \exp (\delta^ {\prime} (a) + \delta^ {\prime \prime} (a)).
\]

因此，a 满足 (3) 当且仅当 b 满足 (4)，定理得证。注。上述证明表明，\(\exp(\hat{\mathbf{L}})\) 是 \(\Delta\) 的子群 \(\Gamma\)，它由满足 (4) 的 b 构成。

因此，\(\Delta\) 的群律可通过指数映射传递到 \(\hat{\mathbf{L}}\)。换言之，\(\hat{\mathbf{L}}\) 是一个完备拓扑群，其复合律 \((a, b) \mapsto a \uplus b\) 由下式给出

\[
a \uplus b = \log (\exp a. \exp b).
\]

如此得到的拓扑群称为豪斯多夫群（由 X 关于 K 导出）。

令 \(g\) 为从自由群 \(\mathbf{F} = \mathrm{F}(\mathbf{X})\) 到 \(\Gamma\) 的同态，满足对 \(g(x) = \exp x\) 有 \(x \in \mathbf{X}\)。由于 \(\exp x - 1 - x = \sum_{n \geqslant 2} x^n / n!\) 的阶 \(\geqslant 2\)，由 §5 第 3 号定理 1，\(g\) 是单射。因此，映射 \(\log \circ g\) 是 \(F\) 到豪斯多夫群的单射同态，并延拓了典范嵌入 \(X \to L\)。

对每个整数 \(m \geqslant 1\)，记 \(\hat{\mathbf{L}}_m\) 为 \(\geqslant m\) 中阶 \(\hat{\mathbf{L}}\) 的元素所成的集合，记 \(\Gamma_m\) 为满足 \(u \in \Gamma\) 的阶 \(u - 1\) 的 \(\geqslant m\) 所成的集合。由第 1 号的注 2，\(\hat{\mathbf{L}}_m = \exp^{-1}(\Gamma_m)\)；由于 \((\Gamma_m)_{m \geqslant 1}\) 是 \(\Gamma\) 上的整中心滤子（§4，第 5 号，命题 2），\((\hat{\mathbf{L}}_m)_{m \geqslant 1}\) 是群 \(\hat{\mathbf{L}}\) 上的整中心滤子。

\subsection{李形式幂级数}

引理 1。设 \(\mathfrak{g}\) 是一个滤李代数（§ 4，第 1 号），\((\mathfrak{g}_{\alpha})_{\alpha \in \mathbf{R}}\) 是其滤子，且 \(\alpha \in \mathbf{R}\)。设 \(\mathbf{P}\) 是 \(n\) 次齐次李多项式，关于不定元 \((\mathbf{T}_i)_{i \in I}\)（§ 2，第 4 号）。则 \(\mathrm{P}((a_i)) \in \mathfrak{g}_{n\alpha}\)，其中 \((a_i)_{i \in I}\) 是由 \(\mathfrak{g}_{\alpha}\) 中元素组成的任意族。

每个次数 \(n \geqslant 2\) 的李多项式都是形如 \([\mathbf{Q}, \mathbf{R}]\) 的项的有限和，其中 \(\mathbf{Q}\) 和 \(\mathbf{R}\) 的次数均为 \(< n\)，且它们的次数之和等于 \(n\)（§ 2，第 7 号，命题 7）。引理由对 \(n\) 作归纳得到。

关于不定元 \((\mathrm{T}_{i})_{i\in\mathbb{I}}\) 的李形式幂级数†（系数在 K 中）是李代数 \(\hat{\mathrm{L}}((\mathrm{T}_{i})_{i\in\mathbb{I}})=\hat{\mathrm{L}}(\mathrm{I})\) 的任一元素。这样的元素 u 可以唯一写成可求和族 \((u_{v})_{v\in\mathbf{N}^{(\mathrm{I})}}\) 的和，其中 \(u_{v}\in\mathrm{L}^{v}(\mathrm{I})\)。

假设 I 是有限集。设 g 是一个完备豪斯多夫滤李代数，满足 \(g = \bigcup_{\alpha > 0} g_{\alpha}\)；令 \(t = (t_i)_{i \in I}\) 为 g 中元素组成的族。

命题 2。满足 \(f_{t} \colon \mathrm{L}(\mathrm{I}) \to \mathfrak{g}\) 的同态 \(f_{t}(\mathrm{T}_{i}) = t_{i}\)（§ 2，第 4 号）可由连续性延拓为唯一的从 \(\hat{f}_{t}\) 到 \(\hat{\mathrm{L}}(\mathrm{I})\) 的连续同态 \(\mathfrak{g}\)。

存在 \(\alpha > 0\)，使得对所有 \(t_i \in \mathfrak{g}_\alpha\) 都有 \(i \in I\)；因此对所有 \(f_t(L^\nu(I)) \subset \mathfrak{g}_{|\nu| \alpha}\)，有 \(\nu\)（引理 1），这蕴含 \(f_t\) 的连续性。

若 \(u \in \widehat{\mathbf{L}}(\mathbf{I})\)，记 \(u((t_i)) = \widehat{f_t}(u)\)。特别地，取 \(\mathfrak{g} = \widehat{\mathbf{L}}(\mathbf{I})\) 时，\(u = u((\mathrm{T}_i))\)；一般情况下，\(u((t_i))\) 称为在李形式幂级数 \(t_i\) 中以 \(T_i\) 代替 \(u((\mathrm{T}_i))\) 所得到的结果。若 \(u = \sum_{v \in \mathbf{N}^{(\mathrm{I})}} u_v\)，其中 \(u_v \in \mathrm{L}^v(\mathrm{I})\)，则族 \((u_v((t_i))_{v \in \mathbf{N}^{(\mathrm{I})}})\) 可求和，并且

\[
u ((t _ {i})) = \sum_ {\nu \in \mathbf {N} ^ {\mathrm{I}}} u _ {\nu} ((t _ {i})).\tag{5}
\]

设 \(\sigma\) 是从 \(\mathfrak{g}\) 到完备豪斯多夫滤李代数 \(\mathfrak{g}'\) 的连续同态，且 \(\mathfrak{g}' = \bigcup_{\alpha > 0} \mathfrak{g}_{\alpha}'\)。对于 \(\boldsymbol{t} = (t_i)_{i \in I}\) 中元素组成的每个有限族 \(\mathfrak{g}\) 以及所有 \(u \in \hat{\mathbf{L}}(\mathbf{I})\)，

\[
\sigma (u ((t _ {i}))) = u ((\sigma (t _ {i}))),\tag{6}
\]

因为同态 \(\sigma \circ f_t\) 是连续的，并且将 \(T_i\) 映到 \(\sigma(t_i)\)（\(i \in I\)）。

令 \(\boldsymbol{u} = (u_j)_{j \in \mathcal{J}}\) 为 \(\hat{\mathrm{L}}(\mathrm{I})\) 中元素组成的有限族，并令 \(v \in \hat{\mathrm{L}}(\mathrm{J})\)；通过用 \(u_j\) 代替 \(T_j\) 并代入 \(v\)，得到元素 \(w = v((u_j)_{j \in \mathcal{J}})\)，它属于 \(\hat{\mathrm{L}}(\mathrm{I})\)，记为 \(v \circ \boldsymbol{u}\)。于是

\[
w ((t _ {i}) _ {i \in I}) = v ((u _ {j} ((t _ {i}) _ {i \in I})) _ {j \in J})\tag{7}
\]

对于 g 中元素组成的每个有限族 \(\boldsymbol{t} = (t_{i})_{i \in I}\) 都成立，这是对等式 \(\hat{f}_{t}\) 施加连续同态 \(w = v((u_{j})_{j \in J})\) 即可看出。

令 \(u = \sum_{\nu \in \mathbb{N}^{\mathbf{I}}} u_{\nu} \in \hat{\mathrm{L}}(\mathrm{I})\)，其中 \(u_{\nu} \in \mathrm{L}^{\nu}(\mathrm{I})\)。从 \(\tilde{u}\colon (t_i)\mapsto u((t_i))\) 到 \(\mathfrak{g}^{\mathbf{I}}\) 的映射 \(\mathfrak{g}\) 是连续的：因为在每个 \(\mathfrak{g}_{\alpha}\) 的开集 \(\alpha > 0\) 中，族 \(\tilde{u}_{\nu}\) 都一致可求和，只需证明每个 \(\tilde{u}_{\nu}\) 连续即可，而这由对 \(|\nu|\) 作归纳立即得到。

\subsection{豪斯多夫级数}

令 \(\{U, V\}\) 为含有两个元素的集合。

定义 1。李代数 \(\mathbf{H} = \mathbf{U} \uplus \mathbf{V} = \log (\exp \mathbf{U} \cdot \exp \mathbf{V})\) 中的元素 \(\hat{\mathbf{L}}_{\mathbf{Q}}(\{\mathbf{U}, \mathbf{V}\})\)（第 2 号）称为关于不定元 \(\mathbf{U}\) 和 \(\mathbf{V}\) 的豪斯多夫级数。

记 \(H_{n}\)（分别记 \(H_{r,s}\)）为 H 的总次数为 n（分别为多重次数 \((r, s)\)）的齐次分量。于是

\[
\mathrm{H} = \sum_ {n \geqslant 0} \mathrm{H} _ {n} = \sum_ {r, s \geqslant 0} \mathrm{H} _ {r, s}, \quad \mathrm{H} _ {n} = \sum_ {\substack {r + s = n \\ r, s \geqslant 0}} \mathrm{H} _ {r, s}.\tag{8}
\]

定理 2。若 \(r\) 和 \(s\) 是满足 \(r + s \geqslant 1\) 的两个正整数，则 \(\mathrm{H}_{r,s} = \mathrm{H}_{r,s}' + \mathrm{H}_{r,s}''\)，其中

\[
(r + s) \mathrm{H} _ {r, s} ^ {\prime} = \tag {9}
\]

\[
\sum_{m\geqslant 1}\frac{(-1)^{m - 1}}{m}\sum_{\substack{r_{1} + \dots +r_{m} = r\\ s_{1} + \dots +s_{m - 1} = s - 1\\ r_{1} + s_{1}\geqslant 1,\ldots ,r_{m - 1} + s_{m - 1}\geqslant 1}}\left(\left(\prod_{i = 1}^{m - 1}\frac{(\mathrm{adU})^{r_{i}}}{r_{i}!}\frac{(\mathrm{adV})^{s_{i}}}{s_{i}!}\right)\frac{(\mathrm{adU})^{r_{m}}}{r_{m}!}\right)(\mathrm{V})
\]

\[
(r + s) \mathrm{H} _ {r, s} ^ {\prime \prime} = \tag {10}
\]

\[
\sum_{m\geqslant 1}\frac{(-1)^{m - 1}}{m}\sum_{\substack{r_{1} + \dots +r_{m - 1} = r - 1\\ s_{1} + \dots +s_{m - 1} = s\\ r_{1} + s_{1}\geqslant 1,\dots ,r_{m - 1} + s_{m - 1}\geqslant 1}}\Bigg(\prod_{i = 1}^{m - 1}\frac{(\operatorname{adU})^{r_{i}}}{r_{i}!}\frac{(\operatorname{adV})^{s_{i}}}{s_{i}!}\Bigg)(U).
\]

在 \(\hat{\mathbf{A}}_{\mathbf{q}}(\{\mathrm{U},\mathrm{V}\})\) 中，exp U .exp V = 1 + W，其中 W = \(\sum_{r+s\geqslant 1}\frac{\mathrm{U}^r}{r!}\frac{\mathrm{V}^s}{s!}\)，从而 H = \(\sum_{m\geqslant 1}(-1)^{m - 1}\mathrm{W}^m /m\)（第 2 号），即：

\[
\mathrm{H}_{r,s} = \sum_{m\geqslant 1}\frac{(-1)^{m - 1}}{m}\sum_{\substack{r_{1} + \dots +r_{m} = r\\ s_{1} + \dots +s_{m} = s\\ r_{1} + s_{1}\geqslant 1,\ldots ,r_{m} + s_{m}\geqslant 1}}\prod_{i = 1}^{m}\frac{\mathrm{U}^{r_{i}}}{r_{i}!}\frac{\mathrm{V}^{s_{i}}}{s_{i}!}.\tag{11}
\]

线性映射 \(\mathbf{P}_n\) 由 \(\mathbf{P}_n(x_1, \ldots, x_n) = \frac{1}{n} \left( \prod_{i=1}^{n-1} (\text{ad } x_i) \right)(x_n)\)（\(n \geqslant 1\) 且 \(x_1, \ldots, x_n\) 属于 \(\{\mathrm{U}, \mathrm{V}\}\)）定义，它是从 \(\mathrm{A}_{\mathbf{Q}}^n(\{\mathrm{U}, \mathrm{V}\})\) 到 \(\mathrm{L}_{\mathbf{Q}}^n(\{\mathrm{U}, \mathrm{V}\})\) 的投影算子（§ 3，第 2 号，命题 1 的推论）；由于 \(\mathrm{H}_{r,s}\) 属于 \(\mathbf{L}_{\mathbf{Q}}^{r+s}(\{\mathrm{U}, \mathrm{V}\})\)，有 \(\mathrm{H}_{r,s} = \mathbf{P}_{r+s}(\mathrm{H}_{r,s})\)。现在

\[
\begin{array}{l} \mathrm{P} _ {r + s} \left(\prod_ {i = 1} ^ {m} \frac {\mathrm{U} ^ {r _ {i}} \mathrm{V} ^ {s _ {i}}}{r _ {i} ! s _ {i} !}\right) \\ = \frac {1}{r + s} \left(\left(\prod_ {i = 1} ^ {m - 1} \frac {(\operatorname{ad} \mathrm{U}) ^ {r _ {i}}}{r _ {i} !} \frac {(\operatorname{ad} \mathrm{V}) ^ {s _ {i}}}{s _ {i} !}\right) \frac {(\operatorname{ad} \mathrm{U}) ^ {r _ {m}}}{r _ {m} !} \frac {(\operatorname{ad} \mathrm{V}) ^ {s _ {m} - 1}}{s _ {m} !}\right) (\mathrm{V}) \end{array}\tag{12}
\]

当 \(s_m \geqslant 1\) 时，以及

\[
\mathrm{P} _ {r + s} \left(\prod_ {i = 1} ^ {m} \frac {\mathrm{U} ^ {r _ {i}}}{r _ {i} !} \frac {\mathrm{V} ^ {s _ {i}}}{s _ {i} !}\right) = \frac {1}{r + s} \left(\left(\prod_ {i = 1} ^ {m - 1} \frac {(\mathrm{adU}) ^ {r _ {i}}}{r _ {i} !} \frac {(\mathrm{adV}) ^ {s _ {i}}}{s _ {i} !}\right) \frac {(\mathrm{adU}) ^ {r _ {m} - 1}}{r _ {m} !}\right) (\mathrm{U})\tag{13}
\]

当 \(r_m \geqslant 1\) 且 \(s_m = 0\) 时。此外，显然，当 \(t\) 时，(ad \(^{p-1}\) ) \(t = 0\) . \(p \geqslant 2\)，且 (ad \(t\) ) \(^0\) . \(t = t\)。因此，当 \(s_m \geqslant 2\) 时，(12) 的两边均为零；当 \(r_m \geqslant 2\) 时，(13) 的两边均为零。定理遂由下述事实得到：\(H_{r,s}'\) 是 (12) 型各项之和，而 \(H_{r,s}''\) 是 (13) 型各项之和。

注。（1）我们已经定义了（§ 3，第 2 号，注）从 A(X) 到 L(X) 的投影算子 Q，它满足对 \(Q(a^{m}) = 0\) 和 \(a \in L(X)\) 有 \(m \geqslant 2\)，且 Q(1) = 0。于是 H = Q(exp H) = Q(exp U.exp V)，立即得到

\[
\mathrm{H} _ {r, s} = \mathrm{Q} \left(\frac {\mathrm{U} ^ {r}}{r !} \frac {\mathrm{V} ^ {s}}{s !}\right) \quad \text { for } r + s \geqslant 1.\tag{14}
\]

\begin{enumerate}
\def\labelenumi{(\arabic{enumi})}
\setcounter{enumi}{1}
\tightlist
\item
  有
\end{enumerate}

\[
\begin{array}{r l} \mathrm{H(U,V)} & \equiv \mathrm{U+V} + \frac {1}{2} [ \mathrm{U,V} ] + \frac {1}{12} [ \mathrm{U,[U,V]} ] \\ & \quad + \frac {1}{12} [ \mathrm{V,[V,U]} ] - \frac {1}{24} [ \mathrm{U,[V,[U,V]} ] ] \end{array} \tag {15}
\]

模 \(\sum_{n\geqslant 5}\mathrm{L}^n (\{\mathrm{U},\mathrm{V}\})\)

\begin{enumerate}
\def\labelenumi{(\arabic{enumi})}
\setcounter{enumi}{2}
\tightlist
\item
  对每个整数 \(\mathrm{H}_{0,n} = \mathrm{H}_{n,0} = 0\)，都有 \(n\neq 1\)，从而
\end{enumerate}

\[
\mathrm{H} (\mathrm{U}, 0) = \mathrm{H} (0, \mathrm{U}) = \mathrm{U}.\tag{16}
\]

另一方面，由于 \([U, -U] = 0\)，

\[
\mathrm{H} (\mathrm{U}, - \mathrm{U}) = 0.\tag{17}
\]

\subsection{豪斯多夫级数中的代入}

由于 K 是含有 Q 的域，豪斯多夫级数可看作系数在 K 中的李形式幂级数。因此，若 g 是满足 \(g = \bigcup_{\alpha > 0} g_{\alpha}\) 的完备豪斯多夫滤李代数，则对于 g 中的 a、b，可以在 H 中以 a、b 代替 U、V（参见第 3 号和 §2 第 5 号的注）。

特别地，令 A 为完备豪斯多夫滤含幺结合代数。对 \(m = \bigcup_{\alpha > 0} A_{\alpha}\)，记 \(m_{\alpha} = A_{\alpha} \cap m\) 和 \(\alpha \in R\)；因此当 \(m_{\alpha} = A_{\alpha}\) 时 \(\alpha > 0\)，而当 \(m_{\alpha} = m\) 时 \(\alpha \leqslant 0\)。赋予 m 括号 \([a, b] = ab - ba\) 后，m 是完备豪斯多夫滤李代数，因而可以应用上述结果。用此记号，得到如下结果，它补充了第 1 号的命题 1。

命题 3。若 \(a \in \mathfrak{m}\)、\(b \in \mathfrak{m}\)，则 \(\exp \mathrm{H}(a, b) = \exp a \cdot \exp b\)。

令 \(a, b\) 属于 \(\mathfrak{m}\)；存在 \(\alpha > 0\)，使得 \(a \in \mathrm{A}_{\alpha}\) 且 \(b \in \mathrm{A}_{\alpha}\)。于是存在连续同态 \(\theta\)，它从 Magnus 代数 \(\hat{\mathbf{A}}(\{\mathbf{U},\mathbf{V}\})\) 到 A，将 U 映到 \(a\)，将 V 映到 \(b\)（ \(\S 5\) ，第 1 号，命题 1）。

将 \(\theta\) 限制在 \(\hat{\mathbf{L}}(\{\mathrm{U},\mathrm{V}\})\) 上，得到从 \(\mathbf{L}(\{\mathrm{U},\mathrm{V}\})\) 的李代数到 \(\mathfrak{m}\) 的连续同态，它将 \(\mathbf{U}\)（分别为 V）映到 \(a\)（分别为 \(b\)）。因此由第 3 号的公式 (6)，\(\theta (\mathrm{H}) = \mathrm{H}(a,b)\)。只需将连续同态 \(\theta\) 施加于关系式两边

\[
\exp \mathrm{H} (\mathrm{U}, \mathrm{V}) = \exp \mathrm{U}. \exp \mathrm{V}
\]

并注意第 1 号的注 3。

注 1。若 \(a\) 与 \(b\) 可交换，则当 \(\mathrm{H}_{r,s}(a,b) = 0\) 时，\(r + s \geqslant 2\)，因为每个次数 \(\geqslant 2\) 的齐次李多项式在 \((a,b)\) 处均为零。于是 \(\mathrm{H}(a,b) = a + b\)，命题 3 恢复了公式

\[
\exp (a + b) = \exp a. \exp b.
\]

命题 4。设 \(\mathfrak{g}\) 是满足 \(\mathfrak{g} = \bigcup_{\alpha >0} \mathfrak{g}_{\alpha}\) 的完备豪斯多夫滤李代数。映射 \((a, b) \mapsto \mathrm{H}(a, b)\) 是 \(\mathfrak{g}\) 上与 \(\mathfrak{g}\) 上的拓扑相容的群律；在该拓扑下，\(0\) 是单位元，且 \(-a\) 是 \(a\) 的逆元，其中对所有 \(a \in \mathfrak{g}\) 均成立。

从 \((a, b) \mapsto \mathrm{H}(a, b)\) 到 \(\mathfrak{g} \times \mathfrak{g}\) 的映射 \(\mathfrak{g}\) 是连续的（第 3 号）；由于映射 \(a \mapsto -a\) 显然连续，只需证明下列关系

\begin{enumerate}
\def\labelenumi{(\arabic{enumi})}
\setcounter{enumi}{17}
\tightlist
\item
\end{enumerate}

\[
\mathrm{H} (\mathrm{H} (a, b), c) = \mathrm{H} (a, \mathrm{H} (b, c))\tag{18}
\]

\[
\mathrm{H} (a, - a) = 0\tag{19}
\]

\[
\mathrm{H} (a, 0) = \mathrm{H} (0, a) = a\tag{20}
\]

对于 \(a, b, c\) 中的 \(\mathfrak{g}\)。由第 3 号公式 (7)，只需在 \(a, b, c\) 为三个不定元且 \(\mathfrak{g} = \hat{\mathbf{L}}\langle \{a, b, c\}\rangle\) 时证明这些公式。现在，指数映射限制在 \(\hat{\mathbf{L}}\langle \{a, b, c\}\rangle\) 上，是到 Magnus 代数 \(\hat{\mathbf{A}}\langle \{a, b, c\}\rangle\) 中的一个嵌入，并且由命题 3：

\[
\exp \mathrm{H} (\mathrm{H} (a, b), c) = \exp \mathrm{H} (a, b). \exp c = \exp a. \exp b. \exp c
\]

\[
\exp \mathrm{H} (a, \mathrm{H} (b, c)) = \exp a. \exp \mathrm{H} (b, c) = \exp a. \exp b. \exp c
\]

\[
\exp \mathrm{H} (a, - a) = \exp a. \exp (- a) = \exp (a - a) = \exp 0
\]

\[
\exp \mathrm{H} (a, 0) = \exp a. \exp 0 = \exp a
\]

\[
\exp \mathrm{H} (0, a) = \exp 0. \exp a = \exp a.
\]

这就证明了关系 (18) 至 (20)。

注。（2）取 g 为李代数 \(\hat{\mathbf{L}}(\mathbf{X})\)。上述命题中引入的群律与第 2 号定义的运算一致。换言之，

\[
a \uplus b = \mathrm{H} (a, b) \quad \text { for   } a, b \text {   in   } \hat {\mathrm{L}} (\mathrm{X});\tag{21}
\]

因此，豪斯多夫群律由豪斯多夫级数给出。

\begin{enumerate}
\def\labelenumi{(\arabic{enumi})}
\setcounter{enumi}{2}
\tightlist
\item
  设 \(\mathfrak{g}\) 是一个带有由下中心列定义的整滤子的李代数 \((\mathcal{C}^m\mathfrak{g})\)。假设存在 \(m \geqslant 1\)，使得 \(\mathcal{C}^m\mathfrak{g} = \{0\}\)。赋予其由滤子 \((\mathcal{C}^m\mathfrak{g})_{n \geqslant 1}\) 导出的拓扑后，\(\mathfrak{g}\) 是豪斯多夫的、完备的，甚至是离散的。于是有 \(\mathrm{P}(a_1, \ldots, a_r) = 0\)，其中 \(a_1, \ldots, a_r\) 属于 \(\mathfrak{g}\)，且这里的 \(\mathrm{P}\) 是次数 \(\geqslant m\) 的任意齐次李多项式；特别地，\(\mathrm{H}_{r,s}(a, b) = 0\) 当 \(r + s \geqslant m\) 时，且级数 \(\mathrm{H}(a, b) = \sum_{r,s} \mathrm{H}_{r,s}(a, b)\) 只有有限个非零项。于是，映射 \((a, b) \mapsto \mathrm{H}(a, b)\) 在 \(\mathfrak{g}\) 上是一个群律，并且是多项式映射（§ 2，第 4 号）。
\end{enumerate}

命题 5。令 \(\mathbf{K}_{r,s}\) 为 \(\mathrm{H}(\mathrm{U} + \mathrm{V}, -\mathrm{U})\) 的多重次数为 \((r,s)\) 的分量。则

\[
\mathrm{K} _ {n, 1} (\mathrm{U}, \mathrm{V}) = \frac {1}{(n + 1) !} (\operatorname{ad} \mathrm{U}) ^ {n} (\mathrm{V}) \quad for n \geqslant 0.
\]

记 \(\mathrm{K}(\mathrm{U},\mathrm{V}) = \mathrm{H}(\mathrm{U} + \mathrm{V}, - \mathrm{U}),\mathrm{K}_1(\mathrm{U},\mathrm{V}) = \sum_{n\geqslant 0}\mathrm{K}_{n,1}(\mathrm{U},\mathrm{V}).\) 记 L（分别为 R）为 \(\hat{\mathbf{A}} (\{\mathbf{U},\mathbf{V}\})\) 上由 U 左乘（分别右乘）的映射

 可写为

\[
\begin{array}{r l} e ^ {\mathrm{U}} \mathrm{V} e ^ {- \mathrm{U}} & = \sum_ {p, q} \frac {\mathrm{U} ^ {p}}{p !} \mathrm{V} \frac {(- \mathrm{U}) ^ {q}}{q !} \\ & = \sum_ {n \geqslant 0} \frac {1}{n !} \left(\sum_ {p + q = n} \frac {n !}{p ! q !} (\mathrm{L} ^ {p} (- \mathrm{R}) ^ {q}). \mathrm{V}\right) \\ & = \sum_ {n \geqslant 0} \frac {1}{n !} (\mathrm{L} - \mathrm{R}) ^ {n}. \mathrm{V} \end{array}
\]

从而

\[
e ^ {\mathrm{U}} \mathrm{V} e ^ {- \mathrm{U}} = \sum_ {n \geqslant 0} \frac {1}{n !} (\mathrm{adU}) ^ {n} \mathrm{V}.\tag{22}
\]

现在模去 \(\sum_{m\geqslant 0}\sum_{n\geqslant 2}\mathrm{A}^{m,n}(\{\mathrm{U},\mathrm{V}\})\) 的理想 \(\mathrm{A}(\{\mathrm{U},\mathrm{V}\})\) 进行计算。对所有 \(n\geqslant 1\)，

\[
(\mathrm{U} + \mathrm{V}) ^ {n} \equiv \mathrm{U} ^ {n} + \sum_ {i = 1} ^ {n - 1} \mathrm{U} ^ {i} \mathrm{V} \mathrm{U} ^ {n - 1 - i}
\]

从而

\[
\begin{array}{r l} (\mathrm{adU}) (\mathrm{U} + \mathrm{V}) ^ {n} & \equiv ((\mathrm{L} - \mathrm{R}) \sum_ {i = 1} ^ {n - 1} \mathrm{L} ^ {i} \mathrm{R} ^ {n - i}). \mathrm{V} \\ & \equiv (\mathrm{L} ^ {n} - \mathrm{R} ^ {n}). \mathrm{V} \\ & \equiv \mathrm{U} ^ {n} \mathrm{V} - \mathrm{VU} ^ {n}. \end{array}
\]

因此

\[
(\mathrm{ad} \mathrm{U}). \mathrm{e} ^ {\mathrm{U} + \mathrm{V}} \equiv e ^ {\mathrm{U}} \mathrm{V} - \mathrm{V} e ^ {\mathrm{U}}\tag{23}
\]

对 n 求和。

另一方面，\(\mathrm{K}_1(\mathrm{U},\mathrm{V})\equiv \mathrm{K}(\mathrm{U},\mathrm{V})\)，且 \(e^{\mathrm{K}_1(\mathrm{U},\mathrm{V})}\equiv 1 + \mathrm{K}_1(\mathrm{U},\mathrm{V})\)，从而

\[
\mathrm{K} _ {1} \equiv e ^ {\mathrm{K}} - 1 \equiv e ^ {\mathrm{U} + \mathrm{V}} e ^ {- \mathrm{U}} - 1
\]

由命题 3 得到

\[
\begin{array}{r l r} (\mathrm{adU}) \mathrm{K} _ {1} & \equiv \mathrm{Ue} ^ {\mathrm{U} + \mathrm{V}} e ^ {- \mathrm{U}} - e ^ {\mathrm{U} + \mathrm{V}} e ^ {- \mathrm{U}} \mathrm{U} \equiv (\mathrm{Ue} ^ {\mathrm{U} + \mathrm{V}} - e ^ {\mathrm{U} + \mathrm{V}} \mathrm{U}) e ^ {- \mathrm{U}} \\ & \equiv (e ^ {\mathrm{U}} \mathrm{V} - \mathrm{Ve} ^ {\mathrm{U}}) e ^ {- \mathrm{U}} & \text {by (23)} \\ & \equiv e ^ {\mathrm{U}} \mathrm{Ve} ^ {- \mathrm{U}} - \mathrm{V} \\ & \equiv \sum_ {n \geqslant 1} \frac {1}{n !} (\mathrm{adU}) ^ {n} \mathrm{V} & \text {by (22)} \\ & \equiv (\mathrm{adU}) \Bigl (\sum_ {n \geqslant 0} \frac {(\mathrm{adU}) ^ {n}}{(n + 1) !} \mathrm{V} \Bigr). \end{array}
\]

于是只需应用 § 2 第 11 号的注。

\section{豪斯多夫级数的收敛性（实数或复数情形）}

本段中假设 K 是带有通常绝对值的域 R 或 C 之一。回忆一下，K 上的可赋范代数是一个（不一定结合的）K 代数，带有拓扑 T，并满足以下性质：

\begin{enumerate}
\def\labelenumi{(\arabic{enumi})}
\item
  \(\mathcal{T}\) 可由范数定义：
\item
  从 \((x,y)\mapsto xy\) 到 \(\mathbf{A}\times \mathbf{A}\) 的映射 \(\mathbf{A}\) 是连续的。
\end{enumerate}

K 上的赋范代数是带有范数的 K 代数 A，满足对 A 中所有 x、y 都有 \(\|xy\|\leqslant\|x\|\|y\|\)。

记 g 为 K 上的完备可赋范李代数。选取 g 上的一个范数和一个数 M > 0，使得

\[
\| [ x, y ] \| \leqslant \mathrm{M} \| x \| \| y \| \quad \text { for } x, y \text { in } \mathfrak {g}.\tag{1}
\]

\subsection{取值于 g 的连续多项式}

设 I 是一个有限集，令 \(\mathbf{P}(\mathfrak{g}^{\mathrm{I}};\mathfrak{g})\)（分别令 \(\hat{\mathbf{P}} (\mathfrak{g}^{\mathrm{I}};\mathfrak{g}))\)）为定义在 \(\mathbf{g}^{\mathrm{I}}\) 上且取值于 \(\mathfrak{g}\) 的连续多项式（分别为带连续分量的形式幂级数）的向量空间。回忆（《微分流形与解析流形》，R，附录），\(\mathbf{P}(\mathfrak{g}^{\mathrm{I}};\mathfrak{g})\) 带有 \(\mathbf{N}^{\mathrm{I}}\) 型分次，而 \(\hat{\mathbf{P}} (\mathfrak{g}^{\mathrm{I}};\mathfrak{g})\) 被认作向量空间 \(\mathbf{P} (\mathfrak{g}^{\mathrm{I}};\mathfrak{g})\) 的完备化，其中的拓扑由与 \(\mathbf{P}(\mathfrak{g}^{\mathrm{I}};\mathfrak{g})\) 的分次相关的滤子定义。此外，\(\mathbf{P}(\mathfrak{g}^{\mathrm{I}};\mathfrak{g})\) 是一个分次李代数，其括号对 \([f,g](\pmb {x}) = [f(\pmb {x}),g(\pmb {x})]\)、\(f, g\) 定义为 \(\mathrm{P}(\mathfrak{g}^{\mathrm{I}}; \mathfrak{g}), x \in \mathfrak{g}^{\mathrm{I}}\)；这个李代数结构可由连续性延拓到 \(\hat{\mathrm{P}}(\mathfrak{g}^{\mathrm{I}}; \mathfrak{g})\)，使其成为完备豪斯多夫滤李代数。

根据 § 6 第 3 号的命题 2，存在唯一的连续李代数同态 \(\phi_{\mathrm{I}}:u\mapsto \tilde{u}\)，它从 \(\hat{\mathbf{L}} (\mathbf{I})\) 到 \(\hat{\mathbf{P}} (\mathbf{g}^{\mathrm{I}};\mathfrak{g})\)，并将指标为 \(i\) 的不定元映到 \(\mathsf{pr}_{\mathrm{i}}\)（对所有 \(i\in \mathbb{I}\)），因为 \(\mathsf{pr}_{\mathrm{i}}\in \mathsf{P}(\mathfrak{g}^{\mathrm{I}};\mathfrak{g})\)。由此，当 \(\tilde{u}\in \mathsf{P}(\mathfrak{g}^{\mathrm{I}};\mathfrak{g})\) 时，\(u\in \mathbf{L}(\mathbf{I})\)；更确切地说，当 \(u\in \mathbf{L}(\mathbf{I})\) 时，\(\tilde{u}\) 正是 § 2 第 4 号中的多项式映射 \((t_i)\mapsto u((t_i))\)。另一方面，显然 \(\phi_{\mathrm{I}}\) 与 \(\mathbf{L}(\mathbf{I})\) 和 \(\mathbf{P}(\mathfrak{g}^{\mathrm{I}};\mathfrak{g})\) 的多重分次相容。若 \(u = \sum_{\nu \in \mathbf{N}^{\mathrm{I}}}u_{\nu}\)，其中对 \(u_{\nu}\in \mathbf{L}^{\nu}(\mathbf{I})\) 有 \(\nu \in \mathbf{N}^{\mathrm{I}}\)，则

\[
\tilde {u} = \sum_ {\nu \in \mathbf {N} ^ {\mathrm{I}}}  \tilde {u} _ {\nu}, \quad \text { where } \tilde {u} _ {\nu} \in \mathrm{P} _ {\nu} (\mathfrak {g} ^ {\mathrm{I}}; \mathfrak {g}).
\]

令 \(\pmb{u} = (u_j)_{j\in J}\) 为 \(\hat{\mathbf{L}} (\mathbf{I})\) 中元素的有限族，令 \(v\in \hat{\mathbf{L}} (\mathbf{J})\)，并令 \(w = v\circ \pmb {u}\)（§ 6，第 3 号）。记 \(\tilde{\pmb{u}} = (\tilde{u}_j)_{j\in \mathbb{J}}\in \mathfrak{J}\)。则

\[
\tilde {v} \circ \tilde {\boldsymbol {u}} = (v \circ \boldsymbol {u}) ^ {\sim}.\tag{2}
\]

这是通过连续延拓 § 6 第 3 号的公式 (7)，并应用《微分流形与解析流形》R 附录第 6 号得到的。

\subsection{由完备赋范李代数定义的群芽}

设 \(H = \sum_{r,s \geqslant 0} H_{r,s} \in \hat{L}(U, V)\) 为豪斯多夫级数（§ 6，第 4 号，定义 1）。我们将证明相应的形式幂级数

\[
\tilde {\mathrm{H}} = \sum_ {r, s \geqslant 0} \tilde {\mathrm{H}} _ {r, s} \in \hat {\mathrm{P}} (\mathfrak {g} \times \mathfrak {g}, \mathfrak {g})\tag{3}
\]

是收敛的（《微分流形与解析流形》R，3.1.1）。

引入下列形式幂级数 \(\eta\in\mathbf{Q}[[U,V]]\)

\begin{enumerate}
\def\labelenumi{(\arabic{enumi})}
\setcounter{enumi}{3}
\tightlist
\item
\item
\end{enumerate}

\[
\begin{array}{l}\eta (\mathrm{U},\mathrm{V}) = -\log (2 - \exp (\mathrm{U} + \mathrm{V}))\\ \qquad = \sum_{m\geqslant 1}\frac{1}{m} (\exp (\mathrm{U} + \mathrm{V}) - 1)^{m}\\ \qquad = \sum_{m\geqslant 1}\frac{1}{m}\sum_{\substack{r_{1},\ldots ,r_{m}\\ s_{1},\ldots ,s_{m}\\ r_{i} + s_{i}\geqslant 1}}\frac{\mathrm{U}^{r_{1}}}{r_{1}!}\frac{\mathrm{V}^{s_{1}}}{s_{1}!}\frac{\mathrm{U}^{r_{2}}}{r_{2}!}\dots \frac{\mathrm{V}^{s_{m}}}{s_{m}!}. \end{array}\tag{6}
\]

因此

\[
\eta (\mathrm{U}, \mathrm{V}) = \sum_ {r, s \geqslant 0} \eta_ {r, s} \mathrm{U} ^ {r} \mathrm{V} ^ {s},\tag{7}
\]

其中

\[
\eta_{r,s} = \sum_{m\geqslant 1}\frac{1}{m}\sum_{\substack{r_{1} + \dots +r_{m} = r\\ s_{1} + \dots +s_{m} = s\\ r_{i} + s_{i}\geqslant 1}}\frac{1}{r_{1}!\ldots r_{m}!s_{1}!\ldots s_{m}!}.\tag{8}
\]

现在令 u 和 v 为满足 \(u + v < \log 2\) 的两个正实数；于是 \(0 \leqslant \exp(u + v) - 1 < 1\)；在 (5) 和 (6) 中分别以 u 代替 U、以 v 代替 V 所得到的级数收敛，且上述计算推出

\[
\sum_ {r, s \geqslant 0} \eta_ {r, s} u ^ {r} v ^ {s} = - \log (2 - \exp (u + v)) <   + \infty .\tag{9}
\]

令 \(r, s \geqslant 0\)，并令 \(\| \tilde{\mathrm{H}}_{r,s} \|\) 为连续多项式 \(\tilde{\mathrm{H}}_{r,s}\) 的范数（《微分流形与解析流形》R，附录，第 2 号）。

引理 1。

\[
\| \widetilde {\mathrm{H}} _ {r, s} \| \leqslant \mathrm{M} ^ {r + s - 1} \eta_ {r, s}.
\]

令 \(r_i, s_i\) 属于 \(\mathbf{N}\)，其中 \(1 \leqslant i \leqslant m\)，且 \(s_m = 1\)；记 \(r = \sum_{i} r_i, s = \sum_{i} s_i\)，并考虑 \(L(\{U, V\})\) 中的下列元素：

\[
Z = \left(\left(\sum_ {i = 1} ^ {m - 1} (\operatorname{ad} U) ^ {r _ {i}} (\operatorname{ad} V) ^ {s _ {i}}\right) (\operatorname{ad} U) ^ {r _ {m}}\right) (V).
\]

于是 \(\tilde{Z} = f \circ p\)，其中 \(f\) 是如下的 \((r + s)\)-线性映射，从 \(\mathfrak{g}^{r + s}\) 到 \(\mathfrak{g}\)：

\[
\begin{array}{l} \left(x _ {1}, \dots , x _ {r}, y _ {1}, \dots , y _ {s}\right) \mapsto \\ \left(\operatorname{ad} \left(x _ {1}\right) \circ \dots \circ \operatorname{ad} \left(x _ {r _ {1}}\right) \circ \operatorname{ad} \left(y _ {1}\right) \circ \dots \circ \operatorname{ad} \left(y _ {s _ {1}}\right) \circ \operatorname{ad} \left(x _ {r _ {1} + 1}\right) \circ \dots \circ \operatorname{ad} \left(x _ {r}\right)\right) \left(y _ {s}\right) \end{array}
\]

其中 \(p\) 是从 \(\mathfrak{g}^2\) 到 \(\mathfrak{g}^{r+s}\) 的如下映射：

\[
(x, y) \mapsto \underbrace {(x , \dots , x} _ {r}, \underbrace {y , \dots , y)} _ {s};
\]

因此 \(\| \tilde{Z}\| \leqslant \| f\| \leqslant M^{r + s - 1}\)（《微分流形与解析流形》R，附录）。将这些不等式应用于 §6 第 4 号公式 (9) 右端的各项，得到：

\[
\| (\mathrm{H} _ {r, s} ^ {\prime}) ^ {\sim} \| \leqslant \frac {\mathrm{M} ^ {r + s - 1}}{r + s} \sum_ {m \geqslant 1} \frac {1}{m} \sum_ {\substack {r _ {1} + \dots + r _ {m} = r \\ s _ {1} + \dots + s _ {m - 1} = s - 1 \\ r _ {1} + s _ {1} \geqslant 1, \dots , r _ {m - 1} + s _ {m - 1} \geqslant 1}} \frac {1}{r _ {1} ! \dots r _ {m} ! s _ {1} ! \dots s _ {m - 1} !}.\tag{10}
\]

类似的论证给出

\[
\begin{array}{l}\| (\mathbf{H}_{r,s}^{\prime \prime})\sim \| \\ \leqslant \frac{\mathbf{M}^{r + s - 1}}{r + s}\sum_{m\geqslant 1}\frac{1}{m}\sum_{\substack{r_{1} + \dots +r_{m - 1} = r - 1\\ s_{1} + \dots +s_{m - 1} = s\\ r_{1} + s_{1}\geqslant 1,\ldots ,r_{m - 1} + s_{m - 1}\geqslant 1}}\frac{1}{r_{1}!\ldots r_{m - 1}!s_{1}!\ldots s_{m - 1}!} \end{array}\tag{11}
\]

从而由 (8)

\[
\left\| \tilde {\mathrm{H}} _ {r, s} \right\| \leqslant   \eta_ {r, s} \frac {\mathrm{M} ^ {r + s - 1}}{r + s} \leqslant \eta_ {r, s} \mathrm{M} ^ {r + s - 1},
\]

这就证明了引理。

命题 1。形式幂级数 \(\tilde{\mathbf{H}}\) 是收敛级数（《微分流形与解析流形》R，3.1.1）；其绝对收敛域（《微分流形与解析流形》R，3.1.4）包含开集

\[
\Omega = \left\{(x, y) \in \mathfrak {g} \times \mathfrak {g} | \| x \| + \| y \| <   \frac {1}{M} \log 2 \right\}.
\]

令 \(u, v\) 为满足 \(> 0\) 的两个 \(u + v < \frac{1}{M} \log 2\) 的实数；则由引理 1

\[
\begin{array}{r l} \mathrm{M} \sum_ {r, s \geqslant 0} \| \tilde {\mathrm{H}} _ {r, s} \| u ^ {r} v ^ {s} \\ & \leqslant \sum_ {r, s \geqslant 0} \eta_ {r, s} \mathrm{M} ^ {r + s} u ^ {r} v ^ {s} = - \log (2 - \exp \mathrm{M} (u + v)) <   + \infty \end{array} \tag {12}
\]

由 (9) 得到。

令 \(h: \Omega \to \mathfrak{g}\) 表示由 \(\tilde{\mathrm{H}}\) 定义的解析函数（《微分流形与解析流形》R，3.2.9），即由公式

\[
h (x, y) = \sum_ {r, s \geqslant 0} \tilde {\mathrm{H}} _ {r, s} (x, y) = \sum_ {r, s \geqslant 0} \mathrm{H} _ {r, s} (x, y) \quad \text { for } (x, y) \in \Omega .\tag{13}
\]

此函数称为关于 M 的 g 的豪斯多夫函数（若不会引起混淆，也简称 g 的豪斯多夫函数）。注意，若 \(\mathrm{H}_{r,s}(\mathrm{U}, -\mathrm{U}) = 0\)，则 \(r + s \geqslant 2\)，从而

\[
h (x, - x) = 0 \quad \text { for } \| x \| <   \frac {1}{2 M} \log 2.\tag{14}
\]

同理

\[
h (0, x) = h (x, 0) = x \quad \text { for } \| x \| <   \frac {1}{\mathrm{M}} \log 2.\tag{15}
\]

命题 2。令

\[
\Omega^ {\prime} = \left\{(x, y, z) \in \mathfrak {g} \times \mathfrak {g} \times \mathfrak {g} | \| x \| + \| y \| + \| z \| <   \frac {1}{M} \log \frac {3}{2} \right\}.
\]

若 \((x,y,z)\in \Omega^{\prime}\)，则

\[
(x, y) \in \Omega , \quad (h (x, y), z) \in \Omega , \quad (y, z) \in \Omega , \quad (x, h (y, z)) \in \Omega\tag{16}
\]

并且

\[
h (h (x, y), z) = h (x, h (y, z)).\tag{17}
\]

令 \((x,y,z)\in \Omega^{\prime}\)；显然 \((x,y)\in \Omega\) 且 \((y,z)\in \Omega\)。此外：

\[
\| h (x, y) \| \leqslant \sum_ {r, s} \| \tilde {\mathrm{H}} _ {r, s} \| \| x \| ^ {r} \| y \| ^ {s},
\]

因而由 (13)

\[
\| h (x, y) \| \leqslant - \frac {1}{M} \log (2 - \exp M (\| x \| + \| y \|)).
\]

现在 \(\mathrm{M}(\|x\| + \|y\|) < \log \frac{3}{2} - \mathrm{M}\|z\|\)；记 \(u = \exp(\mathrm{M}\|z\|)\)；则 \(1 \leqslant u \leqslant \frac{3}{2}\)，并且

\[
\begin{array}{r l} \mathbf {M} (\| h (x, y) \| + \| z \|) <   - \log (2 - \exp (\log \frac {3}{2} - \mathbf {M} \| z \|)) + \mathbf {M} \| z \| \\ = - \log \left(2 - \frac {3}{2 u}\right) + \log u = \log \frac {2 u ^ {2}}{4 u - 3} \\ = \log \left(2 + \frac {2 (u - 1) (u - 3)}{4 u - 3}\right) \leqslant \log 2. \end{array}
\]

同理可见 \((x, h(y, z)) \in \Omega\)。

现在证明 (17)。在李代数 \(\hat{\mathbf{L}}(\{\mathbf{U},\mathbf{V},\mathbf{W}\})\) 中

\[
\mathrm{H} (\mathrm{H} (\mathrm{U}, \mathrm{V}), \mathrm{W}) = \mathrm{H} (\mathrm{U}, \mathrm{H} (\mathrm{V}, \mathrm{W}))
\]

由 §6 第 5 号的命题 4，并由第 1 号公式 (2)，因此在 \(\hat{\mathbf{P}}(\mathfrak{g} \times \mathfrak{g} \times \mathfrak{g}, \mathfrak{g})\) 中有关系

\[
\tilde {\mathrm{H}} \circ (\tilde {\mathrm{H}} \times \mathrm{Id} _ {\mathfrak {g}}) = \tilde {\mathrm{H}} \circ (\mathrm{Id} _ {\mathfrak {g}} \times \tilde {\mathrm{H}}).
\]

由《微分流形与解析流形》R，3.1.9，存在数 \(\varepsilon > 0\)，使得当 \(\|x\|\)、\(\|y\|\) 和 \(\|z\|\) 都 \(\leqslant \varepsilon\) 时公式 (17) 成立。但是，函数 \((x,y,z)\mapsto h(h(x,y),z)\) 和 \((x,y,z)\mapsto h(x,h(y,z))\) 是定义在 \(\Omega'\) 上、取值于 \(\mathfrak{g}\) 的解析函数（《微分流形与解析流形》R，3.2.7）。由于 \(\Omega'\) 连通且它们在 0 的某个邻域内相等，所以它们处处相等（《微分流形与解析流形》R，3.2.5）。

上述结果蕴含：

令 \(\alpha\) 为满足 \(0 < \alpha \leqslant \frac{1}{3M} \log \frac{3}{2}\) 的实数。令

\[
\mathrm{G} = \{x \in \mathfrak {g} \mid \| x \| <   \alpha \},
\]

\(\Theta = \{(x,y) \in G \times G \mid h(x,y) \in G\}\)，并令 \(m: \Theta \to G\) 为 \(h\) 在 \(\Theta\) 上的限制。则：

\begin{enumerate}
\def\labelenumi{(\arabic{enumi})}
\item
  \(\Theta\) 是 \(\mathbf{G} \times \mathbf{G}\) 中的开集，且 \(m\) 是解析的。
\item
  \(x \in G\) 蕴含 \((0, x) \in \Theta\)、\((x, 0) \in \Theta\) 以及 \(m(0, x) = m(x, 0) = x\)。
\item
  \(x \in \mathbf{G}\) 蕴含 \(-x \in \mathbf{G}, (x, -x) \in \Theta, (-x, x) \in \Theta\)，以及
\end{enumerate}

\[
m (x, - x) = m (- x, x) = 0.
\]

\begin{enumerate}
\def\labelenumi{(\arabic{enumi})}
\setcounter{enumi}{3}
\tightlist
\item
  设 \(x, y, z\) 是 \(G\) 中满足 \((x, y) \in \Theta\)、\((m(x, y), z) \in \Theta\)、\((y, z) \in \Theta\) 以及 \((x, m(y, z)) \in \Theta\) 的元素。则 \(m(m(x, y), z) = m(x, m(y, z))\)。
\end{enumerate}

*换言之（第 III 章，§1），若记 \(-x = \sigma(x)\)，则四元组 \((G, 0, \sigma, m)\) 是 \(K_*\) 上的李群芽。

\subsection{完备赋范结合代数中的指数}

在本号中，记 A 为完备赋范含幺结合代数（《一般拓扑学》，第 IX 章，§ 3，第 7 号）。于是对于 A 中的 \(\| x.y\| \leqslant \| x\|\)，有 \(\| y\|\) . \(x,y\)。

设 I 是有限集，令 \(\hat{\mathrm{P}} (\mathrm{A}^{\mathrm{I}};\mathrm{A})\) 为定义在 \(\mathbf{A}^{\mathrm{I}}\) 上、取值于 A 的带连续分量的形式幂级数的向量空间（《微分流形与解析流形》R，附录，第 5 号），其代数结构由下式给出

\[
f. g = m \circ (f, g) \quad \text { for } f, g \text { in } \hat {\mathrm{P}} (\mathrm{A} ^ {\mathrm{I}}; \mathrm{A}),
\]

其中 \(m: \mathbf{A} \times \mathbf{A} \to \mathbf{A}\) 表示 A 上的乘法。仿照第 1 号并使用 § 5 第 1 号的命题 1，定义从 \(u \mapsto \tilde{u}\) 到 \(\hat{\mathrm{A}}(\mathrm{I})\) 的含幺代数连续同态 \(\hat{\mathrm{P}}(\mathrm{A}^{\mathrm{I}}; \mathrm{A})\)，它将指标为 \(i\) 的不定元映到 \(\mathbf{pr}_{\mathrm{i}}\)；这个同态延拓了第 1 号中定义的从 \(\hat{\mathrm{L}}(\mathrm{I})\) 到 \(\hat{\mathrm{P}}(\mathrm{A}^{\mathrm{I}}; \mathrm{A})\) 的李代数同态。若 \(u = \sum_{\nu} u_{\nu}\)，其中对 \(u_{\nu} \in \mathrm{A}^{\nu}(\mathrm{I})\) 有 \(\nu \in \mathbf{N}^{\mathrm{I}}\)，则 \(\tilde{u} = \sum_{\nu} \tilde{u}_{\nu}\)，其中 \(\tilde{u}_{\nu}\) 是多项式映射 \((t_{i})_{i \in I} \mapsto u_{\nu}((t_{i}))\)。

令 \(\boldsymbol{u} = (u_{j})_{j \in I}\) 为 \(\hat{\mathbf{A}}(\mathbf{I})\) 中元素组成的有限族，令 \(v \in \hat{\mathbf{A}}(\mathbf{J})\)，并记 \(w = v \circ \boldsymbol{u}\)（\(\S 5\)，第 1 号）。则

\[
(v \circ \boldsymbol {u}) ^ {\sim} = \tilde {v} \circ \tilde {\boldsymbol {u}}.\tag{18}
\]

这是通过连续延拓 § 5 第 1 号的公式 (2)，并应用《微分流形与解析流形》R 附录第 6 号得到的。

特别地，取 \(\mathrm{I} = \{\mathrm{U}\}\)，把 A 与 \(\mathrm{A}^{(\mathrm{U})}\) 视为同一对象，并考察 \(\tilde{e}\) 中级数 \(\tilde{l}\) 和 \(e(\mathrm{U}) = \sum_{n\geqslant 1}\mathrm{U}^n /n!\) 的像 \(l(\mathrm{U}) = \sum_{n\geqslant 1}(-1)^{n - 1}\mathrm{U}^n /n\) 和 \(\hat{\mathrm{P}} (\mathrm{A};\mathrm{A})\)。于是，对 A 中的 \(\| \tilde{\mathbf{U}}^n\| \leqslant 1\)，若 \(\| x_1\dots x_n\| \leqslant \| x_1\| \dots \| x_n\|\)，则 \(x_{1},\ldots ,x_{n}\)。因此，\(\tilde{e}\)（分别为 \(\tilde{l}\)）的绝对收敛半径为无穷大（分别 \(\geqslant 1\)）。

记 \(e_{A}\)（分别为 \(l_{A}\)）为由收敛级数 \(\tilde{e}\)（分别为 \(\tilde{l}\)）定义的从 A 到 A 的解析映射（分别为从 B 到 A 的解析映射，其中 B 是 A 的开单位球），并记 \(\exp_{\mathbf{A}}(x)=1+e_{\mathbf{A}}(x)\)（对 \(x\in A\)）以及

\[
\log_ {\mathrm{A}} (x) = l _ {\mathrm{A}} (x - 1)
\]

（对 \(x \in \mathbf{A}\)，\(\| x - 1\| < 1\)）。于是

\[
\exp_ {\mathrm{A}} x = \sum_ {n \geqslant 0} \frac {x ^ {n}}{n !} (x \in \mathrm{A})\tag{19}
\]

\[
\log_ {\mathrm{A}} x = \sum_ {n \geqslant 1} (- 1) ^ {n - 1} \frac {(x - 1) ^ {n}}{n} (x \in \mathrm{A}, \| x - 1 \| <   1).\tag{20}
\]

由于 \((e\circ l)(\mathrm{U}) = (l\circ e)(\mathrm{U}) = \mathrm{U}\)（参见 § 6，第 1 号），由 (18) 有 \(\tilde{e}\circ \tilde{l} = \tilde{l}\circ \tilde{e} = \mathrm{Id}_{\mathbf{A}}\)。因此（《微分流形与解析流形》R，3.1.9）

\begin{enumerate}
\def\labelenumi{(\arabic{enumi})}
\setcounter{enumi}{20}
\tightlist
\item
\end{enumerate}

\[
\exp_ {\mathrm{A}} (\log_ {\mathrm{A}} (x)) = x \quad (x \in \mathrm{A}, \| x - 1 \| \leqslant 1)\tag{21}
\]

\[
\log_ {\mathbf {A}} (\exp_ {\mathbf {A}} (x)) = x \quad (x \in \mathrm{A}, \| x \| <   \log 2)\tag{22}
\]

因为当 \(\| x\| < \log 2\) 时有 \(\| \exp_{\mathbf{A}}(x) - 1\| \leqslant \exp \| x\| -1 < 1.\)

最后把 A 看作完备赋范李代数。于是

\[
\| [ x, y ] \| = \| x y - y x \| \leqslant 2 \| x \|. \| y \|.
\]

第 2 号的命题 1 蕴含，形式幂级数 \(\tilde{\mathbf{H}}\) 的绝对收敛域包含集合

\[
\Omega = \{(x, y) \in \mathrm{A} \times \mathrm{A} | \| x \| + \| y \| <   \frac {1}{2} \log 2 \}.
\]

因此 \(\tilde{\mathbf{H}}\) 定义了一个解析函数 \(h\colon \Omega \to \mathrm{A}\)。于是 \(h(x,y) = \sum_{r,s}\mathrm{H}_{r,s}(x,y)\)（参见 § 3，第 1 号，注 4）。

命题 3。当 \(\| x\| + \| y\| < \frac{1}{2}\log 2,\) 时，

\[
\exp_ {A} x. \exp_ {A} y = \exp_ {A} h (x, y).\tag{23}
\]

由 (18) 以及关系 \(e^{\mathrm{U}}e^{\mathrm{V}} = e^{\mathrm{H}(\mathrm{U},\mathrm{V})}\) 可知

\[
m \circ (1 + \tilde {e}, 1 + \tilde {e}) = (1 + \tilde {e}) \circ \tilde {\mathrm{H}}
\]

在 \(\hat{\mathrm{P}} (\mathrm{A}\times \mathrm{A};\mathrm{A})\) 中。因此由《微分流形与解析流形》R，3.1.9，可推出当 \((x,y)\) 充分接近 \((0,0)\) 时 (23) 成立，从而由解析延拓（《微分流形与解析流形》R，3.2.5）得到命题。

\section{豪斯多夫级数的收敛性（超度量情形）}

本段中假设 K 是特征为 0 的非离散完备赋值域，带有超度量绝对值。记 p 为 K 的剩余域的特征（《交换代数》，第 VI 章，§ 3，第 2 号）。

若 \(p \neq 0\)，记 \(a = |p|\)；我们知道（《交换代数》，第 VI 章，§6，第 2、3 号），\(0 < a < 1\)，并且存在唯一的赋值 \(v\)，它定义在 \(\mathbf{K}\) 上、取值于 \(\mathbf{R}\)，其在 \(\mathbf{Q}\) 上的限制是 \(p\)-进赋值 \(v_p\)，并且满足 \(|x| = a^{v(x)}\) 对所有 \(x \in \mathbf{K}\) 成立。此外记：

\[
\theta = \frac {1}{p - 1}.\tag{1}
\]

若 \(p = 0\)，记 \(a\) 为满足 \(0 < a < 1\) 的实数，并记 \(v\) 为定义在 \(\mathbf{K}\) 上、取值于 \(\mathbf{R}\) 的赋值，使得 \(|x| = a^{v(x)}\) 对所有 \(x \in \mathbf{K}\) 成立（同上）。此外，有 \(v(x) = 0\) 对所有 \(x \in \mathbf{Q}^*\) 成立。此外记：

\[
\theta = 0.\tag{2}
\]

\subsection{级数 exp、log 和 H 的 p 进上界}

本号中假设 \(p \neq 0\)。

引理 1。设 n 是一个 \(\geqslant0\) 的整数，且 \(n=n_{0}+n_{1}p+\cdots+n_{k}p^{k}\)，其中 \(0\leqslant n_{i}\leqslant p-1\)，是 n 的 p 进展开。~令 \(\mathrm{S}(n)=n_{0}+n_{1}+\cdots+n_{k}\)。则

\[
v _ {p} (n!) = \frac {n - \mathrm{S} (n)}{p - 1}.\tag{3}
\]

  \(v_{p}(n!) = \sum_{i=1}^{n} v_{p}(i)\)，且整数 \(i\)（介于 1 与 \(n\) 之间）中满足 \(v_{p}(i) \geqslant j\) 的整数个数等于 \([n/p^{j}]\)，即 \(n/p^{j}\) 的整数部分。于是

\[
v _ {p} (n!) = \sum_ {j \geqslant 0} j ([ n / p ^ {j} ] - [ n / p ^ {j + 1} ]) = \sum_ {j \geqslant 1} [ n / p ^ {j} ].
\]

由于 \([n / p^j] = \sum_{i\geqslant j}n_i p^{i - j}\)，引理得证。

引理 2。对每个整数 \(v(n) \leqslant v(n!) \leqslant (n - 1)\theta\)，有 \(v(n) \leqslant (\log n) / (\log p)\) 且 \(n \geqslant 1\)。

由引理 1，\(v(n!) = v_{p}(n!) = (n - S(n))\theta \leqslant (n - 1)\theta\)。

另一方面，\(n \geqslant p^{v(n)}\)，从而 \(v(n) \leqslant (\log n) / (\log p)\)。

令 \(\mathbf{I} = \{\mathbf{U},\mathbf{V}\}\) 为含有两个元素的集合，并令

\[
\mathrm{H} = \sum_ {r, s \geqslant 0} \mathrm{H} _ {r, s} (\mathrm{U}, \mathrm{V}) \in \hat {\mathrm{L}} _ {\mathbf {Q}} (\mathrm{I})
\]

为豪斯多夫级数（§ 6，第 4 号，定义 1）。令 \(\mathbf{Z}_{(p)}\) 为相对于素理想 \(\mathbf{Z}\) 的 \((p)\) 的局部环，令 \((e_b)_{b \in \mathbf{B}}\) 为 \(\mathrm{L}_{\mathbf{Z}_{(p)}}(\mathrm{I})\) 在 \(\mathbf{Z}\) 上的一组基（§ 2，第 11 号，定理 1）。它也是 \(\mathrm{L}_{\mathbf{Q}}(\mathrm{I})\) 在 \(\mathbf{Q}\) 上的一组基。

命题 1。令 \(r\) 和 \(s\) 为两个 \(\geqslant 0\) 的整数。若 \(H_{r,s} = \sum_{b \in B} \lambda_b e_b\)，其中 \(\lambda_b \in \mathbf{Q}\)，是 \(H_{r,s}\) 关于基 \((e_b)_{b \in B}\) 的分解，则

\[
v _ {p} (\lambda_ {b}) \geqslant - (r + s - 1) \theta \quad for all b \in \mathbf {B}.\tag{4}
\]

环 \(\mathrm{A}_{\mathbf{Z}_{(p)}}(\mathrm{I})\) 被认作由词 \(\mathbf{Z}_{(p)}\) 生成的 \(\mathrm{A}_{\mathbf{Q}}(\mathrm{I})\) 的子-\(w \in \mathrm{Mo}(\mathrm{I})\)-模。由于 \(\mathrm{L}_{\mathbf{Z}_{(p)}}(\mathrm{I})\) 是 \(\mathrm{A}_{\mathbf{Z}_{(p)}}(\mathrm{I})\) 的直因子，

\[
\mathrm{L} _ {\mathbf {Z} _ {(p)}} (\mathrm{I}) = \mathrm{A} _ {\mathbf {Z} _ {(p)}} (\mathrm{I}) \cap \mathrm{L} _ {\mathbf {Q}} (\mathrm{I}).\tag{5}
\]

令 \(f\) 为满足 \(f \leqslant (r + s - 1)\theta < f + 1\) 的整数。关系 (4) 等价于对所有 \(v_{p}(\lambda_{b}) \geqslant -f\) 有 \(b \in B\)，即 \(H_{r,s} \in p^{-f}L_{\mathbf{Z}_{(p)}}(I)\)。但由 (5)，这也等价于 \(H_{r,s} \in p^{-f}A_{\mathbf{Z}_{(p)}}(I)\)。

由 §6 第 4 号公式 (11)，只需证明：对每个整数 \(m \geqslant 1\) 以及满足下述条件的所有整数 \(r_1, \ldots, r_m, s_1, \ldots, s_m\)，

\[
r _ {1} + \dots + r _ {m} = r, \qquad s _ {1} + \dots + s _ {m} = s,\tag{6}
\]

\[
r _ {i} + s _ {i} \geqslant 1 \quad \text { for } 1 \leqslant i \leqslant m,
\]

有

\[
v _ {p} (m. r! \dots r _ {m}! s _ {1}! \dots s _ {m}!) \leqslant f.\tag{7}
\]

但是由引理 2，\(v_{p}(r_{i}!s_{i}!) \leqslant (r_{i} + s_{i} - 1)\theta\)，且 \(v_{p}(m) \leqslant v_{p}(m!) \leqslant (m - 1)\theta\)；因此 (7) 的左端有上界

\[
\theta (m - 1 + \sum_ {i = 1} ^ {m} \left(r _ {i} + s _ {i} - 1\right)) = \theta (r + s - 1);
\]

由于它是整数，故它 \(\leqslant f\)，证明完成。

\subsection{赋范李代数}

定义 1。\(\mathbf{K}\) 上的赋范李代数是带有范数的李代数，满足

\begin{enumerate}
\def\labelenumi{(\arabic{enumi})}
\setcounter{enumi}{7}
\tightlist
\item
\end{enumerate}

\[
\| x + y \| \leqslant \sup (\| x \|, \| y \|)\tag{8}
\]

\[
\| [ x, y ] \| \leqslant \| x \|. \| y \|\tag{9}
\]

对 g 中所有 x、y。

在本段其余部分，\(\mathfrak{g}\) 表示完备赋范李代数。

对每个有限集 I，仿照 § 7 第 1 号定义从 \(u \mapsto \tilde{u}\) 到 \(\hat{\mathbf{L}}(\mathbf{I})\) 的连续李代数同态 \(\hat{\mathrm{P}}(\mathfrak{g}^{\mathrm{I}}; \mathfrak{g})\)。如 § 7 中所见，若 \(u = \sum_{\nu} u_{\nu}\)，其中对 \(u_{\nu} \in \mathbf{L}^{\nu}(\mathbf{I})\) 有 \(\nu \in \mathbf{N}^{\mathbf{I}}\)，则 \(\tilde{u} = \sum_{\nu} \tilde{u}_{\nu}\)，其中 \(\tilde{u}_{\nu}\) 是 § 2 第 4 号中定义的多项式映射 \((t_i)_{i \in \mathbf{I}} \mapsto u_{\nu}(t_i)\)。§ 7 第 1 号的复合公式 (2) 仍然成立。

\subsection{由完备赋范李代数定义的群}

设 \(H = \sum_{r,s \geq 0} H_{r,s} \in \hat{L}(\{U, V\})\) 为豪斯多夫级数（§ 6，第 4 号，定义 1）。我们将证明相应的带连续分量的形式幂级数

\[
\tilde {\mathrm{H}} = \sum_ {r, s \geqslant 0} \tilde {\mathrm{H}} _ {r, s} \in \hat {\mathrm{P}} (\mathfrak {g} \times \mathfrak {g}, \mathfrak {g})\tag{10}
\]

是收敛的（《微分流形与解析流形》R，4.1.1）。

令 \(r \geqslant 0\)、\(s \geqslant 0\) 为满足 \(r + s \neq 0\) 的数，并令 \(\| \tilde{\mathrm{H}}_{r,s} \|\) 为连续多项式 \(\tilde{\mathrm{H}}_{r,s}\) 的范数（《微分流形与解析流形》R 附录第 2 号）。

引理 3。

\[
\left\| \tilde {\mathrm{H}} _ {r, s} \right\| \leqslant a ^ {- (r + s - 1) \theta}.
\]

设 B 是关于 I 的 Hall 集，令 \(H_{r,s} = \sum_{b \in B} \lambda_b e_b\) 为 \(H_{r,s}\) 关于 \(L(\{U, V\})\) 中相应基的分解。则

\[
\left| \lambda_ {b} \right| \leqslant a ^ {- (r + s - 1) \theta}.\tag{11}
\]

当 \(p = 0\) 时这显然成立，因为 \(\lambda_b \in \mathbf{Q}\)；当 \(p \neq 0\) 时则由第 1 号命题 1 得到。

此外，

\[
\left\| \tilde {e} _ {b} \right\| \leqslant 1 \quad \text { for } b \in \mathbf {B}.\tag{12}
\]

更一般地，通过对 \(n\) 作归纳可证明：对于关于两个不定元 U 和 V 的每个交错式 \(b\)，其次数为 \(n\)（§2，第 6 号），都有 \(\| \tilde{b} \| \leqslant 1\)。若 \(n = 1\)，则 \(\tilde{b}\) 是从 \(\mathfrak{g} \times \mathfrak{g}\) 到 \(\mathfrak{g}\) 的投影之一，因而范数 \(\leqslant 1\)；若 \(n > 1\)，则存在两个交错式 \(b_1\) 和 \(b_2\)，它们的次数均为 \(< n\)，使得 \(b = [b_1, b_2]\)。由于映射 \(\gamma: (x, y) \mapsto [x, y]\) 是从 \(\mathfrak{g} \times \mathfrak{g}\) 到 \(\mathfrak{g}\) 的范数 \(\leqslant 1\) 的双线性映射，有（《微分流形与解析流形》R，附录，第 4 号）

\[
\| \tilde {b} \| = \| \gamma \circ (\tilde {b} _ {1}, \tilde {b} _ {2}) \| \leqslant \| \tilde {b} _ {1} \|. \| \tilde {b} _ {2} \| \leqslant 1.\tag{13}
\]

关系 (11) 和 (12) 蕴含引理。

命题 2。形式幂级数 \(\tilde{\mathbf{H}}\) 是收敛级数（《微分流形与解析流形》R，4.1.1）。若 \(\mathbf{G}\) 是球 \(\{x\in\mathfrak{g}\mid\|x\|<a^{\theta}\}\)，则 \(\tilde{\mathbf{H}}\) 的绝对收敛域（《微分流形与解析流形》R，4.1.3）包含 \(\mathbf{G}\times\mathbf{G}\)。

若 \(u\) 和 \(v\) 是满足 \(>0\) 且 \(u < a^{\theta}\) 的两个 \(v < a^{\theta}\) 的实数，则由引理 3

\[
\| \tilde {\mathrm{H}} _ {r, s} \| u ^ {r} v ^ {s} \leqslant a ^ {\theta} (u a ^ {- \theta}) ^ {r} (v a ^ {- \theta}) ^ {s}\tag{14}
\]

并且当 \(\| \tilde{\mathrm{H}}_{r,s}\| u^r v^s\) 趋于无穷大时，\(r + s\) 趋于 0。

记 \(h \colon G \times G \to \mathfrak{g}\) 为由 \(\tilde{\mathbf{H}}\) 定义的解析函数（《微分流形与解析流形》R，4.2.4），即由公式

\[
h (x, y) = \sum_ {r, s \geqslant 0} \tilde {\mathrm{H}} _ {r, s} (x, y) = \sum_ {r, s \geqslant 0} \mathrm{H} _ {r, s} (x, y) \quad \text { for } (x, y) \in \mathrm{G} \times \mathrm{G}.\tag{15}
\]

此函数称为 g 的豪斯多夫函数。

令 \((x,y)\in \mathbf{G}\times \mathbf{G}\)。则

\begin{enumerate}
\def\labelenumi{(\arabic{enumi})}
\setcounter{enumi}{15}
\tightlist
\item
\end{enumerate}

\[
\left\| \tilde {\mathrm{H}} _ {r, s} (x, y) \right\| \leqslant \sup (\left\| x \right\|, \left\| y \right\|)\tag{16}
\]

\[
\| h (x, y) \| \leqslant \sup (\| x \|, \| y \|).\tag{17}
\]

\begin{enumerate}
\def\labelenumi{(\arabic{enumi})}
\setcounter{enumi}{16}
\tightlist
\item
  由 (16) 立即得到，而 \(r = s = 0\) 时 (16) 显然成立；若 \(r \geqslant 1\)，则
\end{enumerate}

\[
\begin{array}{r l} \| \tilde {\mathrm{H}} _ {r, s} (x, y) \| & \leqslant \| \tilde {\mathrm{H}} _ {r, s} \| \| x \| ^ {r} \| y \| ^ {s} \\ & \leqslant \| x \| \left(\frac {\| x \|}{a ^ {\theta}}\right) ^ {r - 1} \left(\frac {\| y \|}{a ^ {\theta}}\right) ^ {s} \\ & \leqslant \| x \|; \end{array}
\]

若 \(s \geqslant 1\)，同理可证。

特别地，对于 \(\| h(x,y)\| < a^{\theta}\)，有 \((x,y)\in \mathbf{G}\times \mathbf{G}\)。

命题 3。令 G 为球 \(\{x\in\mathfrak{g}\mid\|x\|<a^{\theta}\}\)。解析映射

\[
h \colon \mathrm{G} \times \mathrm{G} \rightarrow \mathrm{G}
\]

使 G 成为一个群，其中 0 是单位元，并且对所有 \(x \in G\)，-x 是 x 的逆元。此外，若 R 是满足 \(0 < R < a^{\theta}\) 的实数，则球

\[
\{x \in \mathfrak {g} | \| x \| <   \mathrm{R} \}
\]

（分别为 \(\{x \in \mathfrak{g} \mid \|x\| \leqslant R\}\)）是 G 的开子群。

由于 \(\mathrm{H}(\mathrm{U}, -\mathrm{U}) = 0\) 且 \(\mathrm{H}(0, \mathrm{U}) = \mathrm{H}(\mathrm{U}, 0) = \mathrm{U}\)，有 \(h(x, -x) = 0\)，并且

\[
h (0, x) = h (x, 0) = x
\]

对所有 \(x \in G\) 成立。因此还需证明结合律公式

\[
h (h (x, y), z) = h (x, h (y, z)) \quad \text { for } x, y, z \text { in } G.\tag{18}
\]

由于

\[
\mathrm{H} (\mathrm{H} (\mathrm{U}, \mathrm{V}), \mathrm{W}) = \mathrm{H} (\mathrm{U}, \mathrm{H} (\mathrm{V}, \mathrm{W}))
\]

在 \(\hat{\mathbf{L}} (\{\mathrm{U},\mathrm{V},\mathrm{W}\})\) 中（§ 6，第 5 号，命题 4），有

\[
\tilde {\mathrm{H}} \circ (\tilde {\mathrm{H}} \times \mathrm{Id} _ {\mathfrak {g}}) = \tilde {\mathrm{H}} \circ (\mathrm{Id} _ {\mathfrak {g}} \times \tilde {\mathrm{H}})\tag{19}
\]

在 \(\hat{\mathbf{P}} (\mathfrak{g}\times \mathfrak{g}\times \mathfrak{g};\mathfrak{g})\) 中（第 2 号），并且由 (16) 及《微分流形与解析流形》R，4.1.5，(19) 蕴含 (18)。

*换言之（第 III 章，§1），带有豪斯多夫函数的 G 是一个李群。*

\subsection{完备赋范结合代数中的指数}

在本号中，A 表示带有范数 \(x \mapsto \| x \|\) 的含幺结合代数，并满足以下条件：

\[
\begin{array}{c} \| x + y \| \leqslant \sup (\| x \|, \| y \|) \\ \| x y \| \leqslant \| x \|. \| y \| \\ \| 1 \| = 1 \end{array}
\]

对 A 中的 \(x, y\) 成立，并且关于此范数完备。§ 7 第 3 号的第二、三段结果仍然有效。

取 \(\mathbf{I} = \{\mathrm{U}\}\)，并考察 \(\tilde e\) 中级数 \(\tilde{l}\) 和 \(e(\mathrm{U}) = \sum_{n\geqslant 1}\frac{\mathrm{U}^n}{n!}\) 的像 \(l(\mathrm{U}) = \sum_{n\geqslant 1}(-1)^{n - 1}\frac{U^n}{n}\) 和 \(\hat{\mathrm{P}} (\mathrm{A};\mathrm{A})\)。于是：

\begin{enumerate}
\def\labelenumi{(\arabic{enumi})}
\setcounter{enumi}{19}
\tightlist
\item
\end{enumerate}

\[
\left\| \left(\frac {\mathrm{U} ^ {n}}{n !}\right) ^ {\sim} \right\| \leqslant a ^ {- (n - 1) \theta}\tag{20}
\]

\[
\left\| \left(\frac {\mathrm{U} ^ {n}}{n}\right) ^ {\sim} \right\| \leqslant a ^ {- \frac {\log n}{\log p}}\tag{21}
\]

由第 1 号引理 2。于是级数 \(\tilde{e}\)（分别为 \(\tilde{l}\)）的绝对收敛半径为 \(\geqslant a^\theta\)（分别 \(\geqslant 1\)）（《微分流形与解析流形》R，4.1.3）。对 \(R > 0\)，令 \(G_R = \{x \in A \mid \|x\|<R\}\)；记 \(G = G_\theta\)。级数 \(\tilde{e}\)（分别为 \(\tilde{l}\)）定义了解析映射 \(e_{A}\)（分别为 \(l_{A}\)），其定义域分别为 \(G\) 和 \(G_{1}\)，值域为 A。记：

\begin{enumerate}
\def\labelenumi{(\arabic{enumi})}
\setcounter{enumi}{21}
\tightlist
\item
\end{enumerate}

\[
\exp_ {\mathrm{A}} (x) = 1 + e _ {\mathrm{A}} (x) = \sum_ {n \geqslant 0} \frac {x ^ {n}}{n !}\tag{22}
\]

\[
\log_ {\mathrm{A}} (x) = l _ {\mathrm{A}} (x - 1) = \sum_ {n \geqslant 1} (- 1) ^ {n - 1} \frac {(x - 1) ^ {n}}{n} \quad \text { for } x - 1 \in \dot {\mathrm{G}} _ {1}\tag{23}
\]

（若不会引起混淆，省略下标 A。）对 \(x \in G_R\) 和 \(n \geqslant 1\)，

\[
\left\| \frac {x ^ {n}}{n} \right\| \leqslant \left\| \frac {x ^ {n}}{n !} \right\| <   R ^ {n} a ^ {- (n - 1) \theta} = R \left(\frac {R}{a ^ {\theta}}\right) ^ {n - 1}\tag{24}
\]

从而当 \(e_{\mathrm{A}}(\mathrm{G}_{\mathrm{R}}) \subset \mathrm{G}_{\mathrm{R}}, l_{\mathrm{A}}(\mathrm{G}_{\mathrm{R}}) \subset \mathrm{G}_{\mathrm{R}}\) 时，\(R \leqslant a^{\theta}\)。

命题 4。令 R 为满足 \(0 < R \leqslant a^{\theta}\) 的实数。映射 \(\exp_{\mathbf{A}}\) 定义了从 \(\mathbf{G}_{R}\) 到 \(1 + \mathbf{G}_{R}\) 的解析同构，而逆同构是 \(\log_{\mathbf{A}}\) 在 \(1 + \mathbf{G}_{R}\) 上的限制。

\(e(l(\mathbf{X})) = l(e(\mathbf{X})) = \mathbf{X}\)。由 (20)、(21) 以及《微分流形与解析流形》R，4.1.5，得到 \(e_{\mathrm{A}}(l_{\mathrm{A}}(x)) = l_{\mathrm{A}}(e_{\mathrm{A}}(x))\) 对 \(x \in \mathrm{G}_{\mathrm{R}}\) 成立。于是

\[
\begin{array}{l l} \exp_ {\mathrm{A}} (\log_ {\mathrm{A}} x) = x & \text {for} x \in 1 + \mathrm{G} _ {\mathrm{R}} \\ \log_ {\mathrm{A}} (\exp_ {\mathrm{A}} x) = x & \text {for} x \in \mathrm{G} _ {\mathrm{R}} \end{array}
\]

证明完成。

若在 A 上赋予括号 \([x,y]=xy-yx\)，则 A 成为完备赋范李代数，因为 \(\|xy-yx\|\leqslant\sup(\|xy\|,\|yx\|)\leqslant\|x\|. \|y\|\)。第 3 号命题 2 蕴含 \(\tilde{H}\) 的绝对收敛域包含 \(G\times G\)，因此 \(\tilde{H}\) 定义了解析函数 \(h:G\times G\to A\)；于是

\[
h (x, y) = \sum_ {r, s \geqslant 0} \mathrm{H} _ {r, s} (x, y).\tag{25}
\]

命题 5。对于 G 中的 x、y，

\[
\exp_ {\mathrm{A}} x. \exp_ {\mathrm{A}} y = \exp_ {\mathrm{A}} h (x, y).\tag{26}
\]

由 \(e^{\mathrm{U}}e^{\mathrm{V}} = e^{\mathrm{H(U,V)}}\)，从而

\[
m \circ (1 + \tilde {e}, 1 + \tilde {e}) = (1 + \tilde {e}) \circ \tilde {\mathrm{H}}
\]

在 \(\hat{\mathrm{P}}(\mathbf{A} \times \mathbf{A}; \mathbf{A})\) 中（其中 m 表示 A 上的乘法）。于是由命题 2、引理 3 以及《微分流形与解析流形》R，4.1.5，得到命题。


\appendixsection{莫比乌斯函数}

  令 \(n\) 为一个 \(\geqslant 1\) 的整数。若 \(n\) 被某个素数的平方整除，记 \(\mu(n) = 0\)。若 \(n\) 不被任何素数的平方整除，记 \(\mu(n) = (-1)^k\)，其中 \(k\) 是 \(n\) 的素因子个数。如此定义的函数 \(\mu: \mathbf{N}^* \to \{-1, 0, 1\}\) 称为莫比乌斯函数。

回忆一下，给定两个整数 \(n_1 \geqslant 1, n_2 \geqslant 1\)，若 \(n_1|n_2\) 整除 \(n_1\)，就记作 \(n_2\)。

命题。（i）函数 \(\mu\) 是从 \(\mathbf{N}^*\) 到 \(\mathbf{Z}\) 的唯一映射，满足 \(\mu(1) = 1\) 且

\[
\sum_ {d \mid n} \mu (d) = 0\tag{1}
\]

对每个整数 \(n > 1\) 都成立。

\begin{enumerate}
\def\labelenumi{(\roman{enumi})}
\setcounter{enumi}{1}
\tightlist
\item
  设 \(s\) 和 \(t\) 是从 \(\mathbf{N}^*\) 到一个用加法记号表示的交换群的两个映射。为了使
\end{enumerate}

\[
s (n) = \sum_ {d | n} t (d) \quad \text { for   every   integer } n \geqslant 1,\tag{2}
\]

成立，充要条件是

\[
t (n) = \sum_ {d \mid n} \mu (d) s \left(\frac {n}{d}\right) \quad for every integer n \geqslant 1.\tag{3}
\]

（i）中的唯一性断言是显然的，因为 (1) 允许确定 \(\mu(n)\)，可通过对 \(n\) 作归纳完成。下面证明函数 \(\mu\) 满足 (1)。令 \(n\) 为一个 \(>1\) 的整数。令 P 为 \(n\) 的素因子集合，并令 \(n = \prod_{p \in \mathbb{P}} p^{\nu_p(n)}\) 为 \(n\) 的素因子分解。若 \(d\) 是 \(n\) 的一个因子，则 \(\mu(d) = 0\)，除非 \(d\) 形如 \(\prod_{p \in \mathbb{H}} p\)，其中 H 是 P 的一个子集。于是

\[
\begin{array}{r l} \sum_ {d | n} \mu (d) & = \sum_ {\mathrm{H} \subset \mathrm{P}} (- 1) ^ {\mathrm{CardH}} \\ & = \sum_ {k = 0} ^ {\mathrm{CardP}} \left(\frac {\mathrm{CardP}}{k}\right) (- 1) ^ {k} = (1 - 1) ^ {\mathrm{CardP}} = 0. \end{array}
\]


设 \(s\) 和 \(t\) 是从 \(\mathbf{N}^*\) 到一个用加法记号表示的交换群的两个映射。令 \(n \in \mathbf{N}^*\)。若 (2) 成立，则

\[
\begin{array}{r l} \sum_ {d | n} \mu (d) s \left(\frac {n}{d}\right) & = \sum_ {d | n} \mu (d) \sum_ {\delta \mid (n / d)} t (\delta) = \sum_ {d \delta | n} \mu (d) t (\delta) \\ & = \sum_ {\delta | n} t (\delta) \sum_ {d \mid (n / \delta)} \mu (d) = t (n). \end{array}
\]

反之，若 (3) 成立，则

\[
\sum_ {d \mid n} t (d) = \sum_ {d \mid n} \sum_ {\delta \mid d} \mu (\delta) s \left(\frac {d}{\delta}\right) = \sum_ {d \mid n} s (d) \sum_ {\delta \mid (n / d)} \mu (\delta) = s (n),
\]

证明完成。

公式 (3) 称为莫比乌斯反演公式。


\specialchapter{习题}

\exercisesfor{1}

\begin{enumerate}
\def\labelenumi{\arabic{enumi}.}
\item
  假设 K 是特征为 0 的域。令 E 为 K 上有限秩的余交换双代数。证明 \(P(E) = \{0\}\)（应用第 6 号定理 1）。
\item
  令 E 为满足 \(\mathrm{P}(\mathrm{E})=\{0\}\) 的双代数，并令 \((\mathrm{E}_{n})_{n\geqslant0}\) 为与其双代数结构相容的 E 的一个滤子。证明对 n 作归纳可得，对所有 \(E_{n}^{+}=\{0\}\) 有 \(n\geqslant0\)，并由此推出 \(E^{+}=\{0\}\)，即 E 退化为 K。
\item
  令 \(G\) 为幺半群，\(E = K[G]\) 为其代数（《代数》，第 III 章，§2，第 6 号）。
\end{enumerate}

\begin{enumerate}
\def\labelenumi{(\alph{enumi})}
\item
  证明 E 上存在唯一的余代数结构，使得对所有 \(c(g) = g \otimes g\) 有 \(g \in G\)；这个结构与 E 上的代数结构相容，并使 E 成为余交换双代数，其余单位 \(\varepsilon\) 满足对所有 \(\varepsilon(g) = 1\) 有 \(g \in G\)。
\item
  证明 E 的每个本原元素都为零。由此推出，若 \(G \neq \{e\}\)，则双代数 E 不存在与其双代数结构相容的滤子（应用习题 2）。
\item
  假设 \(\mathbf{K}\) 是整环。证明，\(\mathrm{G}\) 中的元素是满足 \(x\in \mathrm{E}\) 且 \(c(x) = x\otimes x\) 的唯一非零元素。证明，\(\mathrm{E}\) 有逆运算 \(i\)（参见《代数》，第 III 章，§ 11，习题 4）当且仅当 \(\mathrm{G}\) 是群，并且在该情形下 \(i(g) = g^{-1}\) 对所有 \(g\in \mathrm{G}\) 成立。
\end{enumerate}

\begin{enumerate}
\def\labelenumi{\arabic{enumi}.}
\setcounter{enumi}{3}
\item
  令 \(\mathbf{E}\) 为余交换双代数，令 \(\mathbf{G}\) 为满足 \(g \in \mathbf{E}\) 且 \(\varepsilon(g) = 1\) 的 \(c(g) = g \otimes g\) 所成的集合。证明 \(\mathbf{G}\) 对乘法封闭，并证明当 \(\mathbf{E}\) 有逆运算时它是一个群。证明，若 \(\mathbf{K}\) 是域，则 \(\mathbf{G}\) 的元素在 \(\mathbf{K}\) 上线性无关。
\item
  令 E 为双代数，令 \(E' = \text{Hom}(E, K)\) 为其对偶，并赋予它由 E 上余代数结构通过对偶性导出的代数结构（《代数》，第 III 章，§11，第 1 号）。令 m 为从 \(u \mapsto u(1)\) 到 \(\mathbf{E}'\) 的同态 \(\mathbf{K}\) 的核；它是 \(\mathbf{E}'\) 的理想。证明，若 \(u \in \mathfrak{m}^2\) 且 \(x \in \mathrm{P}(\mathrm{E})\)，则 \(u(x) = 0\)。当 \(\mathbf{K}\) 是域时，证明 \(\mathrm{P}(\mathrm{E})\) 是 \(\mathfrak{m}^2\) 中 \(\mathbf{E}^+\) 的正交补；推出 \(\dim \mathrm{P}(\mathrm{E}) \leqslant \dim \mathfrak{m} / \mathfrak{m}^2\)，并证明当 \(\mathbf{E}\) 在 \(\mathbf{K}\) 上有限秩时取等号。
\item
  令 E 为双代数，令 \(u \in \mathrm{E}' = \operatorname{Hom}(\mathrm{E}, \mathrm{K})\)（参见习题 5）。令 \(g \in \mathrm{E}\) 满足 \(\varepsilon(g) = 1\) 且 \(c(g) = g \otimes g\)。证明映射 \(u \mapsto u(g)\) 是从 \(\mathrm{E}'\) 到 \(\mathrm{K}\) 的代数同态，且映射 \(y \mapsto gy\) 是余代数 \(\mathrm{E}\) 的自同态；证明当 \(\mathrm{E}\) 有逆运算时，这个自同态是自同构。
\item
  假设 K 是域。令 E 为满足以下条件的双代数：
\end{enumerate}

\begin{enumerate}
\def\labelenumi{(\roman{enumi})}
\item
  E 是余交换的；
\item
  E 有逆运算（《代数》，第 III 章，§ 11，习题 4）；
\item
  \(\mathrm{P(E)} = \{0\}\)；
\item
  E 在 K 上有限秩。
\end{enumerate}

令 \(\mathbf{E}'\) 表示余代数 E 的对偶 K-代数。令 G 表示满足 \(g \in \mathbf{E}\) 且 \(\varepsilon(g) = 1\) 的 \(c(g) = g \otimes g\) 所成的集合；它是一个群（参见习题 4）。

\begin{enumerate}
\def\labelenumi{(\alph{enumi})}
\item
  令 \(g \in \mathbf{G}\)，令 \(\mathfrak{m}_g\) 为从 \(u \mapsto u(g)\) 到 \(\mathbf{E}'\) 的同态 \(\mathbf{K}\) 的核。证明 \(\mathfrak{m}_g = \mathfrak{m}_g^2\)（使用习题 6 将其归结为 \(g = 1\) 的情形；然后应用习题 5）。
\item
  假设 K 代数闭。证明理想 \(m_{g}\) 是 \(E'\) 的唯一极大理想；推出它们的交为 \(\{0\}\)，\(E'\) 可认作乘积 \(K^{G}\)，并且 E 可认作有限群 G 的双代数 K[G]（参见习题 3）。
\end{enumerate}

\begin{enumerate}
\def\labelenumi{\arabic{enumi}.}
\setcounter{enumi}{7}
\item
  证明李代数 \(\mathfrak{g}\) 的包络双代数有逆运算 \(i\)，且 \(i(x) = -x\) 对所有 \(x \in \sigma(\mathfrak{g})\) 成立。
\item
  假设 K 是 Q-代数。令 g 为具有 K-基的李代数，令 U 为其带有典范滤子 \((\mathrm{U}_{n})_{n\geqslant0}\) 的包络双代数。令 \(x\in U^{+}\)，令 n 为一个 \(\geqslant1\) 的整数。证明，x 属于 \(U_{n}^{+}\) 当且仅当 \(c^{+}(x)\) 属于 \(\sum_{\substack{i+j=n\\ i,j\geqslant1}}\mathrm{Im}(\mathrm{U}_{i}^{+}\otimes\mathrm{U}_{j}^{+})\)。
\item
  假设 K 是 Q-代数。令 E 为带有与其双代数结构相容的滤子的余交换 K-双代数。证明态射
\end{enumerate}

\[
f _ {\mathbf {E}} \colon \mathrm{U} (\mathrm{P} (\mathrm{E})) \rightarrow \mathrm{E}
\]

在第 4 号命题 8 中定义的态射是同构，只要假设 P(E) 是自由 K-模（证明与定理 1 相同）。

¶ 11. 设 E 为双代数，I 为有限集。给定一组基 \((e_{\alpha})\)

它由 \(\alpha\) 中的元素 \(\mathbf{N}^{\mathrm{I}}\) 编号，并满足：

\begin{enumerate}
\def\labelenumi{(\roman{enumi})}
\item
  \(e_{0} = 1\)；
\item
  对所有 \(c(e_{\gamma}) = \sum_{\alpha + \beta = \gamma} e_{\alpha} \otimes e_{\beta}\)，有 \(\gamma \in \mathbf{N}^{\mathrm{I}}\)。
\end{enumerate}

条件 (ii) 蕴含，作为余代数，E 同构于 \(\mathrm{TS}(\mathrm{K}^{\mathrm{I}})\) 的对称张量余代数 \(\mathrm{K}^{\mathrm{I}}\)（参见《代数》，第 IV 章，§ 5，第 7 号）。

\begin{enumerate}
\def\labelenumi{(\alph{enumi})}
\tightlist
\item
  令 \(\mathrm{E}' = \mathrm{Hom}(\mathrm{E},\mathrm{K})\) 为 \(\mathbf{E}\) 的对偶代数，令 \(\Lambda = \mathrm{K}[[(\mathrm{X}_i)_{\mathrm{i}\in I}]]\) 为关于由 I 编号的不定元 \(\mathbf{X}_i\) 的形式幂级数代数。对所有 \(u\in \mathrm{E}'\)，令 \(\phi_u\) 为形式幂级数
\end{enumerate}

\[
\sum_ {\alpha} u (e _ {\alpha}) x ^ {\alpha},
\]

其中 \(x^{\alpha} = \prod_{i\in I}\mathrm{X}_{i}^{\alpha (i)}\)

证明 \(u \mapsto \phi_u\) 是从 \(E'\) 到 \(\Lambda\) 的同构，并且这个同构将 \(E'\) 上的逐点收敛拓扑映到 \(\Lambda\) 上的乘积拓扑。

\begin{enumerate}
\def\labelenumi{(\alph{enumi})}
\setcounter{enumi}{1}
\tightlist
\item
  令 \(c_{\alpha \beta \gamma}\) 为代数 \(E\) 关于基 \((e_{\alpha})\) 的结构常数。则
\end{enumerate}

\[
e _ {\alpha} e _ {\beta} = \sum_ {\gamma} c _ {\alpha \beta \gamma} e _ {\gamma}.
\]

用 x 代替 \((\mathbf{X}_{i})_{i\in\mathbb{I}}\)，用 y 代替 \((\mathbf{Y}_{i})_{i\in\mathbb{I}}\)，其中 \(Y_{i}\) 是新的不定元。证明，存在关于变量 x、y 的形式幂级数族 \(f(\boldsymbol{x},\boldsymbol{y})=(f_{i}(\boldsymbol{x},\boldsymbol{y}))_{i\in\mathbb{I}}\)，使得

\[
\boldsymbol {f} (\boldsymbol {x}, \boldsymbol {y}) ^ {\gamma} = \sum_ {\alpha , \beta} c _ {\alpha \beta \gamma} \boldsymbol {x} ^ {\alpha} \boldsymbol {y} ^ {\beta} \quad \text { for   all } \gamma \in \mathbf {N} ^ {\mathrm{I}},
\]

其中如上记 \(x^{\alpha} = \prod_{i} X_{i}^{\alpha(i)}\)，对于 \(y^{\beta}\) 和 \(f(x, y)^{\gamma}\) 也作类似记号。

证明 \(f(x, y)\) 是 \(K\) 上维数为 \(n = \text{Card}(I)\) 的形式群律，其意义见第 I 章 § 1 习题 24†；定义这个形式群律到李代数 \(P(E)\) 的李代数同构，其中 \(e_{\alpha}\)（满足 \(|\alpha| = 1\)）构成一组基。

\begin{enumerate}
\def\labelenumi{(\alph{enumi})}
\setcounter{enumi}{2}
\item
  反之，证明 K 上关于 x、y 的每个形式群律都可由上述步骤得到，并且所得结果在同构意义下唯一。
\item
  当 K 是 Q-代数时，将上述结果应用于在 K 上具有有限基的李代数 g 的包络双代数。推出存在以 g 为李代数的形式群律（且在同构意义下唯一）。
\item
  明确给出与单参数形式群律 \(f(x, y) = x + y\)（分别为 \(f(x, y) = x + y + xy\)）相对应的双代数。
\end{enumerate}

¶ 12. 假设 K 是特征为 p > 0 的域。

\begin{enumerate}
\def\labelenumi{(\alph{enumi})}
\tightlist
\item
  令 \(\mathfrak{g}\) 为李 \(p\)-代数（第 I 章，§ 1，习题 20），令 \(\tilde{\mathbf{U}}\) 为其
\end{enumerate}


限制包络双代数（第 I 章，§2，习题 6）。证明 \(\tilde{\mathbf{U}}\) 上存在唯一的双代数结构，它与其代数结构相容，并且 \(g\) 的元素是本原元。证明 \(\tilde{\mathbf{U}}\) 是余交换的，且 \(\mathrm{P}(\tilde{\mathbf{U}}) = g\)。

\begin{enumerate}
\def\labelenumi{(\alph{enumi})}
\setcounter{enumi}{1}
\item
  若 \(n\) 是一个 \(\geqslant 0\) 的整数，令 \(\tilde{\mathbf{U}}_n\) 表示 \(\tilde{\mathbf{U}}\) 中由乘积 \(x_1\cdots x_n\) 生成的向量子空间，其中 \(x_i \in \mathfrak{g}\) 对所有 \(i\) 成立。证明 \((\tilde{\mathbf{U}}_n)_{n \geqslant 0}\) 是与 \(\tilde{\mathbf{U}}\) 的双代数结构相容的滤子。
\item
  令 \(\mathbf{E}\) 为双代数，令 \(\mathfrak{g} = \mathrm{P}(\mathbf{E})\) 为其本原元素的李代数。映射 \(x\mapsto x^p\) 保持 \(\mathfrak{g}\) 不变，并赋予 \(\mathfrak{g}\) 李 \(p\)-代数结构。证明典范嵌入 \(\mathfrak{g}\to \mathbf{E}\) 可唯一延拓为双代数态射 \(\tilde{\mathbf{U}}\rightarrow \mathbf{E}\)，且该态射是单射（利用以下事实：若 \((x_{i})_{i\in I}\) 是带有全序的 \(\mathfrak{g}\) 的一组基，则单项式 \(\prod_{i\in I}x_i^{n_i}(0\leqslant n_i < p)\) 构成 \(\mathbf{U}\) 的一组基（同上），并仿照引理 2 的证明进行论证）。
\end{enumerate}
